% ============================================================================
% TurboRipassi — Documentazione di progetto
% Parte I:  descrizione in linguaggio naturale (soluzione e obiettivi)
% Parte II: documentazione tecnica (architettura, dati, sicurezza, build)
%
% Compilazione:  pdflatex ripassa-documentazione.tex  (x2 per l'indice)
% ============================================================================
\documentclass[11pt,a4paper]{report}

\usepackage[utf8]{inputenc}
\usepackage[T1]{fontenc}
\usepackage{lmodern}
\usepackage[italian]{babel}
\usepackage[margin=2.5cm]{geometry}
\usepackage{booktabs}
\usepackage{array}
\usepackage{xcolor}
\usepackage{listings}
\usepackage{textcomp}
\usepackage{enumitem}
\usepackage[hidelinks]{hyperref}
\usepackage{parskip}
% I nomi di file e di funzione in \texttt non si spezzano: senza un po' di
% elasticità finiscono per sbordare nel margine.
\setlength{\emergencystretch}{3em}
\usepackage{tikz}
\usetikzlibrary{arrows.meta,positioning,calc,fit,backgrounds}
\usepackage{float}

\definecolor{blu}{HTML}{1E3A8A}
\definecolor{grigiocodice}{HTML}{F3F4F6}
\definecolor{giallo}{HTML}{F59E0B}
\definecolor{grigiobordo}{HTML}{9CA3AF}

% Stili condivisi dai diagrammi TikZ.
\tikzset{
  livello/.style={
    draw=blu, thick, fill=blu!6, rounded corners=3pt,
    minimum width=8.4cm, minimum height=1.05cm, align=center,
    font=\small
  },
  infra/.style={
    draw=giallo!80!black, thick, fill=giallo!12, rounded corners=3pt,
    minimum width=8.4cm, minimum height=0.9cm, align=center, font=\small
  },
  tabella/.style={
    draw=blu, thick, fill=white, rounded corners=2pt,
    minimum width=3.5cm, align=left, font=\scriptsize, inner sep=5pt
  },
  esterno/.style={
    draw=grigiobordo, thick, dashed, fill=gray!5, rounded corners=2pt,
    minimum width=3.2cm, align=center, font=\scriptsize, inner sep=5pt
  },
  freccia/.style={-{Stealth[length=2.4mm]}, thick, draw=blu},
  frecciagrigia/.style={-{Stealth[length=2.4mm]}, thick, draw=grigiobordo},
  etichetta/.style={font=\scriptsize\itshape, text=black!70},
  attore/.style={draw=blu, fill=blu!8, rounded corners=2pt, minimum width=2.2cm,
                 minimum height=0.6cm, font=\scriptsize, align=center},
  vita/.style={draw=grigiobordo, dashed},
  msg/.style={-{Stealth[length=2mm]}, draw=blu},
  msgret/.style={-{Stealth[length=2mm]}, draw=grigiobordo, dashed},
  nota/.style={font=\scriptsize, text=black!65, align=left},
}

\lstset{
  basicstyle=\ttfamily\small,
  backgroundcolor=\color{grigiocodice},
  frame=single, framerule=0pt,
  breaklines=true,
  showstringspaces=false,
  aboveskip=8pt, belowskip=8pt
}

\title{\textbf{\textcolor{blu}{TurboRipassi}}\\[4pt]
\large Documentazione di progetto\\[2pt]
\normalsize App personale per ripassi programmati con ripetizione temporale}
\author{Nikita Piraino \\ \small (codice interamente generato da AI, con supervisione:
vedi il capitolo ``Metodo di lavoro'')}
\date{Revisione 4 --- 15 agosto 2026}

\begin{document}
\maketitle

\vspace{-0.5em}
\begin{center}
\small
\fbox{\begin{minipage}{0.86\textwidth}
\textbf{Stato dei documenti del progetto.} Questo documento descrive lo
\emph{stato attuale} del sistema ed è la fonte autorevole in caso di
divergenza. \texttt{ripassa-app-spec.md} raccoglie i \emph{requisiti
iniziali} (luglio 2026) ed è conservato come documento storico: gli
scostamenti consapevoli rispetto ad esso sono elencati nella sezione
``Scostamenti dalla specifica iniziale''. \texttt{README.md} e
\texttt{BUILD.md} contengono le istruzioni operative.
\end{minipage}}
\end{center}

\tableofcontents

% ============================================================================
\part{La soluzione, in parole semplici}
% ============================================================================

\chapter{Il problema}

Chi studia con il metodo della \emph{ripetizione programmata} rivede lo stesso
argomento a intervalli crescenti: il giorno dopo, dopo una settimana, dopo un
mese, dopo sei mesi. Ogni ``ripasso'' ha bisogno del suo materiale: le foto
degli appunti, i PDF delle dispense, qualche nota scritta al volo.

L'app usata finora per questo scopo (``Genio in 21 Giorni'') ha il concetto
giusto ma un difetto strutturale: su iPad salva i dati su iCloud, su Android
su Google Drive, e i due mondi non si parlano. Il risultato è che il materiale
di studio è frammentato tra dispositivi, e quello che si aggiunge da un
tablet non si ritrova sul telefono.

\chapter{La soluzione: TurboRipassi}

TurboRipassi è un'app personale, pensata per un solo utente, che fa una cosa sola e
la fa su tutti i dispositivi insieme: \textbf{Android (Samsung), iPad e Mac
vedono sempre gli stessi ripassi, le stesse date, gli stessi allegati}, grazie
a un unico account e a un unico archivio centrale (Supabase).

\section{Cosa fa}

\begin{itemize}[nosep]
  \item \textbf{Crea un ripasso} con titolo e note libere.
  \item \textbf{Genera da sola le date di ripasso}: +1 giorno, +1 settimana,
        +1 mese, +6 mesi dal momento della creazione. Un interruttore
        (spento di default) aggiunge anche un ripasso ``+1 ora'' per il
        ripasso immediato.
  \item \textbf{Sposta una data e trascina le successive}: se annoti oggi
        qualcosa studiato due giorni fa, correggi la prima data e l'intera
        scaletta si sposta con lei, invece di restare misurata dal giorno
        sbagliato.
  \item \textbf{Avvisa quando \`e ora}: a ogni data programmata corrisponde un
        promemoria sul dispositivo. Sono notifiche \emph{locali} --- nessun
        servizio push, nessun server che sappia cosa si ripassa.
  \item \textbf{Allega materiale}: foto scattate al momento, immagini dalla
        galleria, PDF e altri file. Gli allegati si possono rinominare,
        riordinare, eliminare.
  \item \textbf{Sincronizza in tempo reale}: una modifica fatta dal telefono
        compare sull'iPad senza fare nulla.
  \item \textbf{Si apre anche senza rete} (in treno, in metro): l'app tiene in
        locale l'elenco dei ripassi, la sessione di accesso e gli allegati dei
        ripassi di \emph{ieri, oggi e domani}, scaricandoli in anticipo e
        cancellando dal dispositivo quelli che non servono pi\`u. In modalit\`a
        aereo si apre su ci\`o che il dispositivo gi\`a sa, dicendo di quando
        \`e la copia che sta mostrando. Il materiale resta comunque sempre al
        sicuro nel cloud.
\end{itemize}

\section{Come si usa}

L'app ha quattro schermate, pi\`u il login:

\begin{enumerate}[nosep]
  \item \textbf{Accesso} — una volta sola per dispositivo, con l'account
        Google oppure con email e password. Da l\`i in poi la sessione resta
        attiva.
  \item \textbf{Home} — l'elenco dei ripassi programmati, una riga per
        data, in due schede: \emph{Ripassi} (oggi e i giorni successivi) e
        \emph{Storico} (i giorni gi\`a passati, con un filtro per vedere solo
        quelli mai completati). Il tondino a sinistra del titolo segna la
        singola data come fatta. In alto la ricerca, il pulsante per
        aggiungerne uno, la striscia ``Come funziona?'' e, a destra, il
        tondino con l'iniziale dell'indirizzo che porta al Profilo.
  \item \textbf{Scheda ripasso} — titolo, note, l'elenco delle date
        programmate (ognuna si pu\`o segnare come completata, posticipare o
        anticipare, trascinando o no le successive) e l'accesso agli allegati.
  \item \textbf{Allegati} — la galleria del materiale: si apre a schermo
        intero, si rinomina, si riordina, si elimina.
  \item \textbf{Profilo} — le cose che si sistemano una volta e poi non si
        guardano pi\`u: metodi di accesso, collegamento a Google Drive, un
        Aiuto a domande e risposte, il pulsante per segnalare un problema e
        l'uscita dall'account. Stavano tutte sulla Home, che \`e la schermata
        che si apre venti volte al giorno --- con ``Esci'' a un tocco
        sbagliato dal titolo.
\end{enumerate}

\section{Le decisioni di fondo}

\begin{itemize}[nosep]
  \item \textbf{Un solo account per tutto}: il problema dell'app precedente
        non deve ripresentarsi. Il difetto di ``Genio in 21 Giorni'' non era
        usare Google Drive, era usare \emph{archivi diversi su piattaforme
        diverse} (iCloud su iPad, Drive su Android). Gli allegati oggi
        stanno su un unico Google Drive, lo stesso su Android, iPad e Mac:
        la scelta e le sue conseguenze sono documentate nel capitolo
        ``Dove vivono gli allegati''.
  \item \textbf{Tutto gratuito} dove possibile; unica eccezione accettabile
        un costo minimo (\textasciitilde 2 \texteuro{}/mese) se davvero necessario.
  \item \textbf{Niente statistiche o grafici}: fuori scopo, per scelta
        esplicita. Anche i promemoria restano \emph{locali}: nessun servizio
        push, nessun token da registrare, nessun server che sappia che cosa e
        quando si ripassa. La differenza \`e nel capitolo dedicato.
  \item \textbf{Codice scritto interamente da AI}, supervisionato ma non
        scritto a mano: lo stack \`e stato scelto anche per essere ben
        documentato e ``generabile'' in modo affidabile.
  \item \textbf{App nativa}, non sito web: serve accesso a fotocamera,
        file e archiviazione locale per l'uso offline.
\end{itemize}

\chapter{Cosa vogliamo realizzare}

\section{Stato attuale (fatto)}

L'MVP \`e completo e funzionante: login (Google e email/password), creazione
ripassi con date automatiche e riprogrammazione a cascata, promemoria locali,
allegati sul Google Drive dell'utente con compressione delle foto,
sincronizzazione in tempo reale, apertura offline (sessione, elenco e allegati
tenuti sul dispositivo), ricerca, storico, schermata di Profilo con Aiuto e
segnalazione problemi. Il codice \`e coperto da una suite di test automatici
(logica di dominio, repository, hook del Controller e tutte e cinque le
schermate; conteggio e copertura nel cap.~\ref{sec:test}) e ha superato un
audit di sicurezza; \texttt{lint}, typecheck e
test girano a ogni push in integrazione continua. Dettagli nella parte
tecnica.

\emph{Nota sull'esecuzione:} l'autorizzazione a Google Drive richiede uno
schema di redirect derivato dall'\texttt{applicationId}, che Expo Go non pu\`o
onorare. Lo sviluppo quotidiano resta possibile con Expo Go per tutto il
resto, ma \textbf{gli allegati funzionano solo su una build installata}
(development build o APK): \`e la ragione principale per cui la ``Fase 1'' non
\`e pi\`u opzionale.

\`E stata inoltre predisposta la rete di sicurezza necessaria per distribuire
l'app a persone diverse dallo sviluppatore: se qualcosa va storto sul
dispositivo di un utente, l'errore viene segnalato automaticamente (Sentry) e
l'app mostra una schermata di errore con un pulsante ``Riprova'' invece di
chiudersi o restare bianca. Il servizio va attivato inserendo una chiave di
configurazione al momento della distribuzione: senza, l'app funziona
comunque, semplicemente non segnala nulla.

\section{Prossimi passi}

\begin{enumerate}[nosep]
  \item \textbf{App installate stabilmente} (``Fase 1''): l'APK per il Samsung
        \`e costruito nel cloud di Expo (profilo \texttt{preview}, gratis) e
        installato --- 27,7~MB, dopo il lavoro sul peso descritto nel capitolo
        sulla build. Resta l'app per iPad, firmata dal Mac con l'Apple ID
        gratuito. Con l'app installata il login persiste e non serve pi\`u Expo
        Go n\'e la rete del Mac. La procedura completa \`e in \texttt{BUILD.md}.
  \item \textbf{Validazione sul Mac M2}: verificare che l'app per iPad giri
        nativamente sul Mac come app ``Designed for iPad'' (rischio
        principale individuato dalla specifica).
  \item \textbf{Decisione consapevole sul rinnovo settimanale}: con l'Apple
        ID gratuito la firma dell'app iPad scade ogni 7 giorni e va
        rinnovata ricollegando l'iPad al Mac. L'alternativa \`e l'Apple
        Developer Program (99~\$/anno). Si parte gratis e si valuta
        sull'uso reale.
  \item \textbf{Attivare il crash reporting}: creare il progetto su
        \texttt{sentry.io} e registrare la chiave
        (\texttt{EXPO\_PUBLIC\_SENTRY\_DSN}) su EAS --- oggi
        \texttt{app.json} porta ancora un segnaposto. Il codice \`e gi\`a
        pronto: manca solo la chiave. Passaggio necessario prima di dare
        l'app a utenti reali, e ora vale doppio: senza DSN restano invisibili
        non solo i crash, ma anche le segnalazioni che gli utenti scrivono di
        loro iniziativa dal Profilo --- il pulsante c'\`e, risponde
        onestamente ``non attive in questa build'', e non manda niente a
        nessuno.
\end{enumerate}

\section{Possibili sviluppi futuri (oggi fuori scopo)}

Non fanno parte dell'MVP, ma l'architettura non li preclude: scrittura
offline (creare ripassi senza rete, con coda di sincronizzazione), uso
multi-utente (le protezioni dati per-utente sono gi\`a attive sul server, in
previsione di proporre l'app a terzi), pubblicazione sugli store.

% ============================================================================
\part{Documentazione tecnica}
% ============================================================================

\chapter{Stack e vincoli}

\begin{center}
\begin{tabular}{@{}lll@{}}
\toprule
\textbf{Livello} & \textbf{Scelta} & \textbf{Motivo} \\
\midrule
Frontend & React Native + Expo SDK 54 & un codebase per Android/iPadOS/macOS \\
Dati e sync & Supabase (Postgres, & gratuito per uso personale, realtime \\
            & Realtime, Auth) & nativo, un solo account \\
Allegati (binari) & Google Drive dell'utente & niente tetto da 1~GB; vedi cap.~\ref{sec:drive} \\
                  & (scope \texttt{drive.file}) & \\
Cache locale & expo-sqlite + FileSystem & lettura offline senza servizi esterni \\
Promemoria & expo-notifications & solo notifiche locali: nessun servizio \\
           & (sole notifiche locali) & push, nessun token, nessun server \\
Autenticazione & Supabase Auth & Google OAuth (PKCE) + email/password \\
Crash reporting & Sentry (\texttt{@sentry/} & visibilit\`a sui crash delle build \\
                & \texttt{react-native}) & distribuite; disattivabile senza DSN \\
Test & jest-expo + & logica, repository e hook testabili \\
     & Testing Library RN & senza dispositivo \\
Qualit\`a & ESLint + tsc strict + CI & confini di architettura verificati \\
\bottomrule
\end{tabular}
\end{center}

Limiti del piano gratuito Supabase (verificati, luglio 2026): 500~MB database,
1~GB storage file, 5~GB/mese di banda in uscita, pausa automatica del progetto
dopo 7 giorni senza richieste. Il tetto di 1~GB \`e stato \emph{il} vincolo
dimensionante del progetto ed \`e la ragione per cui i binari degli allegati
non stanno pi\`u su Supabase (capitolo~6). Su Postgres restano solo i
metadati, che a 500~MB non sono un problema realistico. Le immagini vengono
comunque compresse prima dell'upload (ridimensionamento a 1600px, JPEG 70\%),
perch\'e lo spazio consumato ora \`e quello del Drive personale dell'utente.

\chapter{Architettura}

Separazione \textbf{Model / Controller / View} obbligatoria, pi\`u un livello
trasversale \texttt{config} che nella specifica iniziale non era previsto e che
il codice ha reso necessario: l'infrastruttura (client Supabase, SDK Sentry,
storage cifrato, lettura della configurazione) non appartiene al dominio, ma
serve a tutti e tre i livelli.

\begin{figure}[H]
\centering
\resizebox{\textwidth}{!}{%
\begin{tikzpicture}[node distance=0.55cm]
  \node[livello] (view) {\textbf{View} --- \texttt{src/view}\\
    \scriptsize schermate e componenti React Native; tema isolato in \texttt{theme.ts}\\
    \scriptsize nessuna logica di dominio, nessun accesso ai dati};
  \node[livello, below=of view] (ctrl) {\textbf{Controller} --- \texttt{src/controller}\\
    \scriptsize custom hook: stato, orchestrazione, politica di retry, messaggi d'errore\\
    \scriptsize nessun JSX};
  \node[livello, below=of ctrl] (model) {\textbf{Model} --- \texttt{src/model}\\
    \scriptsize tipi di dominio, logica pura, repository (Supabase / Drive), cache SQLite\\
    \scriptsize non conosce n\'e la UI n\'e React};
  \node[infra, below=0.75cm of model] (config) {\textbf{config} --- \texttt{src/config}
    \; \scriptsize(livello trasversale: client Supabase, Sentry, SecureStore, env)};

  \draw[freccia] (view) -- node[etichetta, right=1pt] {hook: stato + azioni} (ctrl);
  \draw[freccia] (ctrl) -- node[etichetta, right=1pt] {repository, funzioni pure} (model);

  \draw[frecciagrigia] ($(view.west)+(0,0)$) to[out=180,in=180,looseness=1.6] (config.west);
  \draw[frecciagrigia] ($(model.east)+(0,0)$) to[out=0,in=0,looseness=1.6] (config.east);
\end{tikzpicture}}
\caption{I livelli e le sole direzioni di dipendenza consentite. Le frecce
grigie indicano che ogni livello pu\`o importare \texttt{config}, mentre
\texttt{config} non importa nessuno di essi.}
\end{figure}

\section{Le regole, e come sono verificate}

\begin{enumerate}[nosep]
  \item Nessuna View importa il client Supabase, un repository, la cache o il
        client Drive: ogni accesso ai dati passa da un hook del Controller.
        Restano leciti dalla View i \emph{tipi} di dominio e le funzioni pure
        (\texttt{ripassiLogic}, \texttt{fileUtils}), che non sono I/O.
  \item Il Model non importa n\'e View n\'e Controller: le dipendenze puntano
        solo verso il basso.
  \item Il Controller non importa la View.
  \item \texttt{config} non importa nessun livello applicativo. Questa regola
        non \`e teorica: il gestore dei token Drive stava in \texttt{config} e
        importava dal Model, creando un ciclo; \`e stato spostato in
        \texttt{model/drive/}.
\end{enumerate}

\textbf{Le quattro regole sono eseguibili.} Fino alla revisione precedente
vivevano solo in questa pagina, ed erano gi\`a state violate una volta e
ripristinate a mano durante l'audit. Ora sono scritte in
\texttt{eslint.config.js} come \texttt{no-restricted-imports} per zona, con un
messaggio che spiega la regola violata: \texttt{npm run lint} fallisce, in
locale e in CI. Una regola di architettura che nessuno pu\`o infrangere per
distrazione vale pi\`u di una regola scritta bene.

\section{Contratti fra i livelli}

Ogni livello espone un'interfaccia dichiarata, non inferita dall'implementazione.

\begin{center}
\small
\begin{tabular}{@{}p{3.1cm}p{11.4cm}@{}}
\toprule
\textbf{Confine} & \textbf{Contratto} \\
\midrule
Model $\rightarrow$ Controller &
  \texttt{RipassiRepo}, \texttt{AllegatiRepo}, \texttt{DriveClient},
  \texttt{DriveTokenManager}: interfacce con un'implementazione di default.
  Gli hook le ricevono come parametro con valore predefinito, quindi un test
  passa un oggetto finto senza toccare il sistema dei moduli. \\
\addlinespace
Controller $\rightarrow$ View &
  \texttt{StatoAuth}, \texttt{StatoRipassi}, \texttt{StatoCache},
  \texttt{ContestoRipassi}: interfacce esportate e annotate come tipo di
  ritorno degli hook, cos\`i una divergenza fallisce nel file dell'hook e non
  nelle cinque schermate che lo consumano. \\
\bottomrule
\end{tabular}
\end{center}

\subsection*{Gestione degli errori: chi lancia, chi traduce, chi parla}

\`E una politica di livello, uniforme in tutto il codice:

\begin{itemize}[nosep]
  \item Il \textbf{Model} lancia l'errore originale, senza tradurlo e senza
        registrarlo: \texttt{if (error) throw error}. Non sa chi lo user\`a.
  \item Il \textbf{Controller} decide. Traduce con \texttt{traduciErrore}
        (Model), sceglie se ritentare con \texttt{conRetry} (Model) e segnala
        con \texttt{mostraErrore}, unico canale verso l'utente, che a sua volta
        registra l'evento su Sentry.
  \item La \textbf{View} non gestisce errori: mostra la stringa che riceve.
\end{itemize}

\begin{figure}[H]
\centering
\resizebox{\textwidth}{!}{%
\begin{tikzpicture}[node distance=0.42cm and 1.1cm]
  \node[esterno] (pg) {Postgres / Drive\\\texttt{42501}, \texttt{403}, \dots};
  \node[tabella, right=of pg, minimum width=3.1cm] (repo) {\textbf{Model} --- repository\\\scriptsize rilancia invariato};
  \node[tabella, right=of repo, minimum width=3.4cm] (hook) {\textbf{Controller} --- hook\\\scriptsize \texttt{conRetry} se transitorio\\\scriptsize \texttt{traduciErrore}};
  \node[tabella, right=of hook, minimum width=2.9cm] (ui) {\textbf{View}\\\scriptsize mostra e basta};
  \node[esterno, below=1.05cm of hook] (sentry) {Sentry\\\scriptsize (via \texttt{mostraErrore})};

  \draw[freccia] (pg) -- (repo);
  \draw[freccia] (repo) -- node[etichetta, above] {throw} (hook);
  \draw[freccia] (hook) -- node[etichetta, above] {italiano} (ui);
  \draw[frecciagrigia] (hook) -- (sentry);
\end{tikzpicture}}
\caption{Percorso di un errore tecnico fino all'utente. Il testo originale non
supera mai il Controller, tranne nell'unico caso documentato al
capitolo~\ref{sec:errori}.}
\end{figure}

\section{Mappa del codice}

Il codice \`e organizzato per \emph{livello} e, dentro ogni livello, per
\emph{dominio} (\texttt{ripassi}, \texttt{allegati}, \texttt{drive},
\texttt{cache}, \texttt{auth}), cos\`i che una modifica funzionale tocchi una
cartella e non tre.

\begin{center}
\scriptsize
\begin{tabular}{@{}p{5.9cm}p{8.9cm}@{}}
\toprule
\textbf{File} & \textbf{Responsabilit\`a} \\
\midrule
\multicolumn{2}{@{}l}{\textbf{Radice}}\\
\texttt{App.tsx} & gating autenticazione, navigazione stack, ErrorBoundary, init crash reporting \\
\midrule
\multicolumn{2}{@{}l}{\textbf{config/} --- infrastruttura trasversale}\\
\texttt{env.ts} & unico lettore di configurazione: env $\rightarrow$ \texttt{app.json/extra}, riconoscimento placeholder \\
\texttt{supabase.ts} & unica istanza del client; PKCE; refresh token via AppState \\
\texttt{secureAuthStorage.ts} & sessione cifrata in Keychain/Keystore (a blocchi da 1800 byte) \\
\texttt{crashReporting.ts} & unico punto di accesso a Sentry: init, DSN, \texttt{reportError}, \texttt{inviaSegnalazione} \\
\texttt{driveConfig.ts} & client ID per piattaforma, scope \texttt{drive.file}, cartella applicativa \\
\texttt{notifications.ts} & impianto delle notifiche locali: permesso, canale Android, primo piano \\
\midrule
\multicolumn{2}{@{}l}{\textbf{model/} --- dominio, dati, I/O}\\
\texttt{types.ts} & tipi del dominio (specchiano lo schema Postgres) \\
\texttt{ripassi/occorrenzeDates.ts} & calcolo puro delle date (+1h/+1g/+1s/+1m/+6m), comparatore, \texttt{ricalcolaSuccessive} \\
\texttt{ripassi/ripassiLogic.ts} & classificazione per giornata (ripassi/storico), ordinamento totale, ricerca \\
\texttt{ripassi/ripassiRepo.ts} & \texttt{RipassiRepo}: CRUD ripassi/occorrenze su Supabase, spostamento transazionale \\
\texttt{ripassi/ripassiOffline.ts} & istantanea dell'elenco sul dispositivo (un JSON, versionato) \\
\texttt{notifiche/notificheLogic.ts} & logica pura: quali occorrenze meritano un promemoria, e il diff \\
\texttt{notifiche/notificheRepo.ts} & \texttt{NotificheRepo}: pianifica e cancella, con l'id dell'occorrenza \\
\texttt{allegati/allegatiRepo.ts} & \texttt{AllegatiRepo}: upload su Drive + metadati, riordino, eliminazione \\
\texttt{drive/driveTypes.ts} & interfacce ed errori tipizzati del backend Drive \\
\texttt{drive/driveRepo.ts} & client REST Drive: cartella, upload in due passi, download, delete \\
\texttt{drive/driveAuth.ts} & OAuth Drive con PKCE + refresh token, persistiti in SecureStore \\
\texttt{cache/cacheLogic.ts} & logica pura della rotazione (finestra, selezione, eliminazione) \\
\texttt{cache/localCache.ts} & I/O cache: SQLite \texttt{cache\_allegati} + filesystem \\
\texttt{auth/oauthRedirect.ts} & parsing e confronto degli URL di redirect \\
\texttt{auth/codiciUsati.ts} & codici OAuth monouso gi\`a spesi, condivisi dai due flussi \\
\texttt{auth/sessioneLocale.ts} & copia della sessione in SecureStore, leggibile senza rete \\
\texttt{shared/errorMessages.ts} & errori tecnici $\rightarrow$ messaggi italiani, per categoria \\
\texttt{shared/retry.ts} & backoff esponenziale sui soli errori transitori \\
\texttt{shared/fileUtils.ts} & estensione da nome/mime, riconoscimento immagini \\
\texttt{shared/account.ts} & garantisce che l'accesso sia agganciato a un account \\
\midrule
\multicolumn{2}{@{}l}{\textbf{controller/} --- stato e orchestrazione}\\
\texttt{AuthContext.tsx} & istanza unica di \texttt{useAuth} per tutta l'app \\
\texttt{RipassiContext.tsx} & istanza unica di ripassi + cache (una sola subscription Realtime) \\
\texttt{avvisoErrore.ts} & \texttt{mostraErrore}: traduce, registra su Sentry, mostra. Unico canale \\
\texttt{useConnettivita.ts} & stato online/offline (NetInfo) per la View \\
\texttt{useLocalCache.ts} & rotazione all'apertura (max 1/giorno) ed esito dell'ultima \\
\texttt{useRitento.ts} & ``sto ritentando'', detto alle schermate: contatore sul gancio di \texttt{conRetry} \\
\texttt{useNotificheRipassi.ts} & tiene i promemoria del dispositivo allineati all'elenco \\
\texttt{auth/useAuth.ts} & sessione, login Google (PKCE) / email, logout; compone useDriveAuth \\
\texttt{auth/useDriveAuth.ts} & autorizzazione Drive: unico proprietario dello stato \\
\texttt{auth/oauthLogin.ts} & helper senza stato del login Google (redirect, messaggi, attese) \\
\texttt{auth/useAccountDrive.ts} & account e cartella Drive, caricati su richiesta \\
\texttt{ripassi/useRipassi.ts} & stato ripassi, CRUD, politica di retry, Realtime \\
\texttt{ripassi/useFormRipasso.ts} & stato del form e orchestrazione del salvataggio \\
\texttt{allegati/useAllegati.ts} & upload a lotti, rinomina, riordino, eliminazione, apertura \\
\texttt{allegati/fileDispositivo.ts} & picker nativi e apertura file locali (senza stato React) \\
\midrule
\multicolumn{2}{@{}l}{\textbf{view/} --- interfaccia}\\
\texttt{theme/theme.ts} & palette blu/giallo, spaziature, raggi: isolata \\
\texttt{lib/format.ts} & date in italiano, etichette Oggi/Domani/Ieri, et\`a della copia locale \\
\texttt{lib/calendarUtils.ts} & griglia del mese (puro, settimana da luned\`i) \\
\texttt{screens/*} & Login, Home, FormRipasso, DettaglioAllegati, Profilo \\
\texttt{components/ui.tsx} & Button, Card, Badge, SectionTitle, Casella, Voce, TestataATendina \\
\texttt{components/Aiuto.tsx} & domande e risposte a tendina, testo statico \\
\texttt{components/SegnalaProblema.tsx} & form di segnalazione con il contesto dell'ultimo errore \\
\texttt{components/allegati.tsx} & miniatura, lista allegati, visualizzatore a schermo intero \\
\texttt{components/occorrenze.tsx} & righe delle date: memorizzata e in anteprima \\
\texttt{components/PulsantiAllegato.tsx} & i tre pulsanti (foto, galleria, file), condivisi \\
\texttt{components/OccorrenzaEditor.tsx} & modale con calendario per spostare/completare \\
\texttt{components/PannelloDrive.tsx} & quale account e quale cartella ricevono gli allegati \\
\texttt{components/PannelloAccount.tsx} & metodi di accesso e collegamento di Google \\
\texttt{components/ErrorBoundary.tsx} & cattura errori di render, segnala e mostra il fallback \\
\midrule
\multicolumn{2}{@{}l}{\textbf{Fuori da \texttt{src}}}\\
\texttt{supabase/schema.sql} & tabelle, trigger, RLS, Realtime, funzioni di riordino e spostamento (idempotente) \\
\texttt{supabase/migrations/} & modifiche allo schema, da eseguire in ordine \\
\texttt{docs/account-identita.md} & perché i ripassi seguono la persona e non il metodo di accesso \\
\texttt{plugins/withAbiRelease.js} & config plugin: una sola architettura nativa per APK \\
\texttt{scripts/autolinkingPerProfilo.js} & esclude la catena del dev client dall'autolinking, fuori dal profilo \texttt{development} \\
\texttt{scripts/releaseLocale.js} & build di release locale con lo stesso peso di quella EAS \\
\texttt{eslint.config.js} & regole Expo + confini fra livelli \\
\texttt{.github/workflows/ci.yml} & lint, typecheck, test e copertura a ogni push \\
\texttt{src/**/\_\_tests\_\_/} & suite Jest (conteggio e copertura nel cap.~\ref{sec:test}) \\
\texttt{supabase/tests/} & test SQL su RLS e modello account/identit\`a \\
\bottomrule
\end{tabular}
\end{center}

\chapter{Modello dati}

Tre tabelle di contenuto su Postgres (Supabase), tutte con RLS attiva, pi\`u
due tabelle di identit\`a e tre depositi solo locali sul dispositivo (una
tabella SQLite e due file). I binari degli allegati non sono in nessuna di
esse: stanno su Google Drive, e \texttt{allegati.storage\_path} ne conserva
l'identificativo.

I dati appartengono all'\texttt{account} — la persona — e non alla riga di
\texttt{auth.users}, che rappresenta un singolo \emph{modo di accedere}. Un
account ha N \texttt{identita}: email e password, Google, e quante se ne
aggiungeranno. \`E ci\`o che fa s\`i che la stessa persona ritrovi gli stessi
ripassi qualunque pulsante di login prema.

\begin{figure}[H]
\centering
\resizebox{\textwidth}{!}{%
\begin{tikzpicture}[node distance=1.0cm and 1.5cm]
  \node[tabella] (account) {%
    \textbf{account} \textit{(la persona)}\\[2pt]
    \underline{id} : uuid\\
    email\_canonica : text\\
    email\_verificata : bool};

  \node[tabella, left=1.6cm of account] (identita) {%
    \textbf{identita} \textit{(un modo di entrare)}\\[2pt]
    \underline{auth\_user\_id} : uuid $\rightarrow$ auth.users\\
    account\_id : uuid\\
    provider : text \textit{(email, google)}\\
    email\_verificata : bool};

  \node[tabella, below=1.1cm of account] (ripassi) {%
    \textbf{ripassi}\\[2pt]
    \underline{id} : uuid\\
    account\_id : uuid $\rightarrow$ account\\
    user\_id : uuid \textit{null (audit)}\\
    titolo : text\\
    note : text \textit{null}\\
    created\_at / updated\_at};

  \node[tabella, below left=1.1cm and -0.6cm of ripassi] (occorrenze) {%
    \textbf{occorrenze}\\[2pt]
    \underline{id} : uuid\\
    ripasso\_id : uuid\\
    account\_id : uuid\\
    scheduled\_at : timestamptz\\
    is\_manual\_1h : bool\\
    is\_completed : bool};

  \node[tabella, below right=1.1cm and -0.6cm of ripassi] (allegati) {%
    \textbf{allegati}\\[2pt]
    \underline{id} : uuid\\
    ripasso\_id : uuid\\
    account\_id : uuid\\
    display\_name / original\_file\_name\\
    \textbf{storage\_path} : text \textit{(ID Drive)}\\
    order\_index : int\\
    mime\_type / size\_bytes};

  \node[esterno, right=1.4cm of allegati] (drive) {Google Drive\\cartella\\\texttt{ripassiProgrammati}};

  \node[tabella, below=1.15cm of occorrenze, draw=grigiobordo] (cache) {%
    \textbf{cache\_allegati} \textit{(solo SQLite locale)}\\[2pt]
    \underline{allegato\_id} : text\\
    local\_uri : text\\
    cached\_at : text \textit{(ultimo giorno in finestra)}};

  \draw[freccia] (identita) -- node[etichetta, above] {N\,:\,1} (account);
  \draw[freccia] (ripassi) -- node[etichetta, right=1pt, pos=0.55] {N\,:\,1 \; cascade} (account);
  \draw[freccia] (occorrenze) -- node[etichetta, left=1pt, pos=0.55] {N\,:\,1 \; cascade} (ripassi);
  \draw[freccia] (allegati) -- node[etichetta, right=1pt, pos=0.55] {N\,:\,1 \; cascade} (ripassi);
  \draw[frecciagrigia] (allegati) -- node[etichetta, above] {1\,:\,1} (drive);
  \draw[frecciagrigia] (cache.west) to[out=180,in=180,looseness=1.2]
        node[etichetta, left] {copia locale} ($(allegati.west)+(0,-0.3)$);
\end{tikzpicture}}
\caption{Le tabelle Postgres (RLS attiva su tutte), la tabella di sola cache
sul dispositivo, e il legame con il file su Drive. Il tratteggio segna ci\`o
che non vive su Postgres. Da notare la direzione delle frecce in alto: i
ripassi pendono dall'\texttt{account}, non dall'identit\`a con cui si \`e
entrati.}
\end{figure}

\section{\texttt{ripassi}}
\begin{center}
\begin{tabular}{@{}llp{8.6cm}@{}}
\toprule
\textbf{Campo} & \textbf{Tipo} & \textbf{Note} \\
\midrule
id & uuid & PK, \texttt{gen\_random\_uuid()} \\
account\_id & uuid & FK \texttt{account}: il proprietario. Default \texttt{account\_corrente()} \\
user\_id & uuid & FK \texttt{auth.users} \textit{null}, default \texttt{auth.uid()}: solo audit \\
titolo & text & \\
note & text & nullable, testo semplice \\
created\_at / updated\_at & timestamptz & \texttt{updated\_at} mantenuto da trigger \\
\bottomrule
\end{tabular}
\end{center}

\section{\texttt{occorrenze}}
\begin{center}
\begin{tabular}{@{}llp{8.6cm}@{}}
\toprule
\textbf{Campo} & \textbf{Tipo} & \textbf{Note} \\
\midrule
id & uuid & PK \\
ripasso\_id & uuid & FK $\rightarrow$ ripassi, \texttt{ON DELETE CASCADE} \\
account\_id & uuid & proprietario; ridondante ma necessario per RLS efficiente \\
user\_id & uuid & \textit{null}, solo audit \\
scheduled\_at & timestamptz & \\
is\_manual\_1h & boolean & true solo per l'occorrenza +1 ora \\
is\_completed & boolean & default false \\
created\_at / updated\_at & timestamptz & \\
\bottomrule
\end{tabular}
\end{center}

Alla creazione di un ripasso vengono inserite le occorrenze a
\textbf{+1 giorno, +1 settimana, +1 mese, +6 mesi} (\texttt{occorrenzeDates.ts});
l'occorrenza \textbf{+1 ora} solo se l'interruttore \`e attivo. Nota
documentata: il 31 gennaio +1 mese trabocca su marzo (comportamento standard
di JavaScript, coperto da test).

\subsection*{Spostare una data trascina le successive}

Annoti oggi qualcosa che hai studiato due giorni fa: correggevi la prima data e
le altre quattro restavano misurate dal giorno sbagliato, con l'intera scaletta
silenziosamente sfasata. Quello che la spaziatura misura \`e la distanza dallo
studio, quindi una prima data corretta all'indietro deve portarsi dietro le
altre.

\texttt{ricalcolaSuccessive} (Model, funzione pura) traduce lo spostamento in
una \textbf{traslazione rigida} delle date seguenti, non in un ricalcolo dagli
offset. Le due cose coincidono finch\'e nessuno ha toccato niente, e dove
divergono \`e la traslazione a essere corretta: una data che l'utente aveva
gi\`a trascinato su una domenica libera si tiene la sua domenica invece di
essere sovrascritta da \texttt{base~+~1m}. Non serve nemmeno sapere da quale
offset venisse una riga --- informazione che lo schema, per inciso, non
registra.

Restano ferme le occorrenze precedenti a quella spostata, quelle gi\`a
completate (che tu abbia ripassato in un certo giorno \`e un fatto sul passato,
non una previsione da correggere) e quelle con una data illeggibile, saltate
invece di essere riscritte come \texttt{Invalid Date}.

La scrittura \`e una sola chiamata transazionale (\texttt{sposta\_occorrenze},
migrazione 0003), per la stessa ragione del riordino degli allegati: $N$ update
separati lascerebbero la scaletta spostata a met\`a se la connessione cade nel
mezzo, e su mobile cade. Le date scritte sono assolute, quindi il retry resta
sicuro. La spunta ``sposta anche le successive'' \`e attiva di default ---
correggere una data e lasciare ferme le altre non \`e quasi mai ci\`o che si
intendeva --- e sparisce quando non ci sarebbe nulla da trascinare.

\section{\texttt{allegati}}
\begin{center}
\begin{tabular}{@{}lll@{}}
\toprule
\textbf{Campo} & \textbf{Tipo} & \textbf{Note} \\
\midrule
id & uuid & PK \\
ripasso\_id & uuid & FK $\rightarrow$ ripassi, cascade \\
account\_id & uuid & proprietario, per RLS \\
user\_id & uuid & \textit{null}, solo audit \\
display\_name & text & modificabile dall'utente \\
original\_file\_name & text & nome originale sul dispositivo \\
storage\_path & text & ID del file su Google Drive (vedi cap.~\ref{sec:drive}) \\
order\_index & integer & riordino manuale \\
mime\_type / size\_bytes & text / bigint & \\
\bottomrule
\end{tabular}
\end{center}

\section{\texttt{cache\_allegati} (solo SQLite on-device)}
\begin{center}
\begin{tabular}{@{}lll@{}}
\toprule
\textbf{Campo} & \textbf{Tipo} & \textbf{Note} \\
\midrule
allegato\_id & TEXT & PK \\
local\_uri & TEXT & percorso nel filesystem del dispositivo \\
cached\_at & TEXT & ultimo giorno (locale) in cui era in finestra \\
\bottomrule
\end{tabular}
\end{center}

\section{Le due copie locali (file, non tabelle)}
\label{sec:copielocali}

Accanto alla cache degli allegati vivono altri due depositi sul dispositivo,
entrambi nati per la stessa ragione --- far aprire l'app quando la rete non
c'\`e --- ed entrambi cancellati al logout, perch\'e ci\`o che contengono \`e
dell'utente e non del dispositivo.

\begin{center}
\small
\begin{tabular}{@{}p{3.6cm}p{3.0cm}p{7.9cm}@{}}
\toprule
\textbf{Deposito} & \textbf{Dove} & \textbf{Che cosa contiene, e perch\'e l\`i} \\
\midrule
sessione locale &
  SecureStore, chiave \texttt{turboripassi.} \texttt{sessione-locale} &
  L'ultima sessione vista, riletta senza toccare la rete. \`E cifrata come
  l'originale perch\'e porta lo stesso refresh token: in chiaro sul filesystem
  sarebbe un posto peggiore. \\
\addlinespace
istantanea dell'elenco &
  \texttt{ripassi-} \texttt{offline.json}, nella document directory &
  Ripassi, occorrenze e metadati degli allegati come li ha restituiti
  l'ultimo caricamento riuscito, pi\`u l'istante in cui il server li ha
  confermati. Un file e non una tabella: si legge e si riscrive intero, una
  volta per caricamento, e SQLite aggiungerebbe uno schema da migrare per
  nessuna query che serva. Un campo \texttt{versione} fa scartare
  un'istantanea scritta da una forma precedente invece di migrarla. \\
\bottomrule
\end{tabular}
\end{center}

Entrambi i moduli \textbf{ingoiano i propri errori}. \`E deliberato e vale la
pena dirlo: esistono per far aprire l'app quando qualcosa \`e andato storto, ed
essere loro la ragione per cui non si apre annullerebbe il motivo per cui ci
sono. Un disco pieno non deve trasformare una schermata che funziona in un
errore; il caricamento successivo riscrive.

\paragraph{Nota su \texttt{cached\_at}.} Il campo \`e rimasto, ma il suo
significato \`e cambiato e non \`e pi\`u un criterio di decisione: la
rotazione elimina per \emph{appartenenza alla finestra} (vedi cap.~\ref{sec:cache}), non per
data di download. Oggi \texttt{cached\_at} \`e solo un'informazione
diagnostica: ``l'ultimo giorno in cui questo allegato serviva''. Il criterio
originale, ancora descritto nella specifica iniziale, cancellava e
riscaricava di continuo file perfettamente validi.

\chapter{Sincronizzazione cross-device}

\begin{itemize}[nosep]
  \item Ogni scrittura passa da Supabase; ogni dispositivo mantiene
        \textbf{una sola} subscription Realtime (in \texttt{RipassiContext})
        sulle tre tabelle: alla ricezione di un evento il Controller ricarica
        e la View si ri-renderizza.
  \item Gli eventi ravvicinati collassano in un solo ricaricamento (finestra
        di 200~ms). Senza, ogni mutazione ne produceva due --- quello
        esplicito e l'eco del proprio evento Realtime --- e un caricamento a
        lotti di $N$ allegati ne produceva $N+1$, ciascuno una rilettura
        completa di ripassi, occorrenze e allegati con conseguente
        ri-render della lista, proprio mentre l'app era occupata a
        comprimere e caricare foto.
  \item Un contatore di sequenza scarta le risposte fuori ordine: con pi\`u
        ricaricamenti in volo la risposta pi\`u vecchia poteva arrivare per
        ultima e rimettere in lista uno stato gi\`a superato. Si manifesta
        quando la rete \`e lenta, cio\`e quando gli eventi si accodano.
  \item Conflitti: \emph{last-write-wins} su \texttt{updated\_at} (trigger
        server-side). Sufficiente per un singolo utente su pi\`u dispositivi.
  \item Assunzione dichiarata: creare/modificare richiede connessione; le
        copie locali servono alla \emph{lettura} offline (cap.~\ref{sec:offline}).
        La scrittura offline (outbox) \`e fuori MVP.
\end{itemize}

\chapter{Cache locale e rotazione}
\label{sec:cache}

Finestra mantenuta in locale: \textbf{ieri, oggi, domani} — calcolata sui
\emph{giorni locali del dispositivo}, non UTC (un ripasso all'1:00 di notte
conta per il giorno giusto).

Ad ogni apertura dell'app (al massimo una volta al giorno per sessione):
\begin{enumerate}[nosep]
  \item si scaricano gli allegati dei ripassi con almeno un'occorrenza in
        finestra (se gi\`a presenti si aggiorna solo \texttt{cached\_at});
  \item si eliminano file e righe di cache degli allegati \textbf{non pi\`u
        appartenenti} alla finestra corrente — inclusi quelli eliminati da
        remoto.
\end{enumerate}

\begin{figure}[H]
\centering
\resizebox{\textwidth}{!}{%
\begin{tikzpicture}[node distance=0.8cm]
  \node[tabella, minimum width=3.6cm] (a) {\textbf{1. Finestra}\\\scriptsize ripassi con almeno\\\scriptsize un'occorrenza in\\\scriptsize [ieri, oggi, domani]\\\scriptsize $\rightarrow$ i loro allegati};
  \node[tabella, right=of a, minimum width=3.6cm] (b) {\textbf{2. Per ciascuno}\\\scriptsize gi\`a in cache e il file\\\scriptsize esiste? aggiorna e basta\\\scriptsize altrimenti scarica\\\scriptsize (senza retry)};
  \node[tabella, right=of b, minimum width=3.6cm] (c) {\textbf{3. Pulizia}\\\scriptsize righe il cui allegato\\\scriptsize NON \`e pi\`u in finestra\\\scriptsize $\rightarrow$ elimina file e riga};
  \node[esterno, below=0.9cm of b] (esito) {\textbf{esito}: disponibili / falliti\\\scriptsize banner in Home + evento Sentry\\\scriptsize (solo se non \`e assenza di rete)};

  \draw[freccia] (a) -- (b);
  \draw[freccia] (b) -- (c);
  \draw[frecciagrigia] (b) -- (esito);
\end{tikzpicture}}
\caption{La rotazione, al massimo una volta al giorno per sessione. Il dato
remoto su Drive non viene mai toccato.}
\end{figure}

Decisione rilevante emersa dall'audit: l'eliminazione avviene per
\emph{appartenenza alla finestra} e non per data di download
(\texttt{cached\_at}); il criterio originale cancellava e riscaricava di
continuo i file ancora validi ma scaricati due o pi\`u giorni prima
(regressione coperta da test). Il dato remoto non viene mai toccato dalla
rotazione. Al logout la cache viene svuotata (i file appartengono all'utente).

\section{Un fallimento silenzioso non \`e pi\`u silenzioso}

Un singolo download che fallisce non deve fermare gli altri, quindi il ciclo
li racchiude uno per uno. Fino alla revisione precedente per\`o l'errore
veniva semplicemente scartato: nessun messaggio, nessun conteggio, nessun
evento su Sentry. La conseguenza era la peggiore possibile per questa
funzionalit\`a --- token Drive revocato, file cancellato dall'utente dal
proprio Drive, spazio esaurito sul dispositivo: in tutti questi casi la cache
non si riempiva mai e \textbf{nessuno lo sapeva}. L'utente lo scopriva in
treno, senza rete, davanti a un allegato che non si apriva; e il crash
reporting, integrato proprio per ``sapere che qualcosa \`e andato storto'',
non riceveva nulla perch\'e nessun crash era avvenuto.

Ora la rotazione restituisce un esito (\texttt{disponibili}, \texttt{falliti},
\texttt{perRete}): la Home mostra una riga quando parte del materiale non \`e
disponibile offline, e i fallimenti \emph{anomali} generano un evento Sentry
per rotazione (uno, non uno per file). L'assenza di rete \`e esclusa di
proposito: \`e un esito previsto, e segnalarla riempirebbe la dashboard di
eventi su cui non si pu\`o agire.

\paragraph{Una rotazione senza rete non consuma la giornata.} Il terzo campo,
\texttt{perRete}, non \`e un dettaglio diagnostico: decide se la rotazione vada
rifatta. Il limite di ``una volta al giorno'' presuppone che il tentativo
\emph{sia avvenuto}, e un download fallito perch\'e la connessione non c'era non
\`e avvenuto affatto. Il difetto \`e emerso quando l'app ha iniziato ad aprirsi
offline (cap.~\ref{sec:offline}): prima, senza rete, non si arrivava nemmeno
alla rotazione. Ora s\`i --- si apre, la rotazione parte contro una rete
assente, fallisce, e con la vecchia regola avrebbe marcato la giornata come
spesa: gli allegati dei ripassi di domani non sarebbero stati scaricati per il
resto della giornata, per quanto presto fosse tornata la linea. Quando almeno
un fallimento \`e di rete la prenotazione viene rilasciata, e il primo
cambiamento della lista --- che \`e esattamente ci\`o che una riconnessione
produce --- riprova.

\chapter{Aprire senza rete}
\label{sec:offline}

La cache degli allegati (cap.~\ref{sec:cache}) esisteva dalla specifica, e la Home prometteva
da sempre ``sei offline --- vedi i ripassi gi\`a scaricati''. La promessa era
falsa, e in un modo che si notava solo nel momento in cui contava: i \emph{file}
erano sul dispositivo, ma i ripassi e le occorrenze a cui appartengono venivano
richiesti a Supabase a ogni apertura. Senza rete la promessa si risolveva in una
schermata vuota con sopra un errore, e gli allegati scaricati erano irraggiungibili
per costruzione --- non restava nessuna riga da cui aprirli.

Il difetto era in realt\`a doppio, e la met\`a pi\`u grave veniva prima: l'app
non arrivava nemmeno a mostrare la lista, perch\'e si fermava al login.

\section{Il primo ostacolo: \texttt{getSession()} non \`e una lettura locale}

Supabase la sessione la conserva gi\`a (\texttt{secureAuthStorage}), ma la
restituisce solo attraverso \texttt{getSession()}, che locale non \`e: quando
l'access token \`e scaduto --- e dopo un'ora lo \`e sempre --- prova prima a
rinnovarlo. Senza connessione quel tentativo ritenta con backoff per una
trentina di secondi e poi riporta \emph{nessuna sessione}. L'app non era rotta e
l'utente non era stato disconnesso; semplicemente, da ci\`o che
\texttt{getSession()} restituisce non c'era modo di dirlo. Il risultato era un
avvio a freddo in modalit\`a aereo che mostrava una rotella lunghissima e poi la
schermata di login --- l'unica schermata che, senza rete, non serve a niente.

La risposta mancante \`e \texttt{sessioneLocale.ts}: una copia della sessione
scritta ogni volta che ne arriva una, rileggibile in millisecondi senza toccare
la rete.

\begin{figure}[H]
\centering
\resizebox{\textwidth}{!}{%
\begin{tikzpicture}[node distance=0.7cm]
  \node[tabella, minimum width=3.9cm] (a) {\textbf{1. Il dispositivo}\\\scriptsize copia in SecureStore\\\scriptsize letta senza rete\\\scriptsize $\rightarrow$ decide il primo frame\\\scriptsize e chiude \texttt{loading}};
  \node[tabella, right=of a, minimum width=4.1cm] (b) {\textbf{2. Il server}\\\scriptsize \texttt{getSession()}, quando\\\scriptsize arriva --- o non arriva};
  \node[tabella, right=of b, minimum width=4.6cm] (c) {\textbf{3. Che cosa se ne fa}\\\scriptsize sessione $\rightarrow$ sostituisce\\\scriptsize uscita vera $\rightarrow$ cancella\\\scriptsize \textbf{irraggiungibile} $\rightarrow$ non tocca};

  \draw[freccia] (a) -- (b);
  \draw[freccia] (b) -- (c);
\end{tikzpicture}}
\caption{L'ordine \`e tutto il punto: il disco decide che cosa si vede subito,
la rete pu\`o solo migliorarlo. Il terzo caso della terza colonna \`e la ragione
per cui esiste l'intera forma.}
\end{figure}

La distinzione che tiene in piedi il meccanismo \`e fra ``nessuna sessione'' e
``nessuna risposta''. Le due cose raggiungono l'hook nello stesso modo --- un
\texttt{null}, da \texttt{getSession()} o dall'evento \texttt{INITIAL\_SESSION}
--- e offline il \texttt{null} \`e l'esito \emph{normale}. Trattarlo come un
logout \`e ci\`o che riportava al login. Ora solo \texttt{SIGNED\_OUT} \`e una
decisione; un \texttt{null} accompagnato da un errore di rete, con una copia sul
dispositivo, non cambia niente. Un flag (\texttt{sessioneDaDispositivo}) ricorda
che quella a schermo \`e una copia e non una conferma, e si spegne appena arriva
l'una o l'altra --- una arriva sempre, perch\'e in React Native il ticker del
token continua a girare e il primo refresh che passa chiude la questione.

\section{Il secondo: l'elenco non lo teneva nessuno}

\texttt{ripassiOffline.ts} scrive l'istantanea dopo ogni caricamento riuscito e
la rilegge all'avvio (cap.~\ref{sec:copielocali}). Il punto delicato \`e che le
due letture --- disco e server --- partono \emph{insieme}, non in sequenza: con
una connessione che funziona il server vince quasi sempre la corsa e
dell'istantanea non si accorge nessuno, e quando non la vince, una lista di
stamattina \`e una prima schermata incomparabilmente migliore di una vuota. A
rendere sicura la corsa \`e una sola guardia: l'istantanea viene mostrata
\emph{solo finch\'e} il server non ha ancora risposto, quindi un disco lento non
pu\`o sovrascrivere righe pi\`u fresche.

\section{Dirlo, e dire di quando}

Una lista salvata \`e da leggere se \`e di stamattina e da mettere in dubbio se
\`e della settimana scorsa, e ``salvati'' da solo non dice n\'e l'uno n\'e
l'altro: il banner della Home riporta l'et\`a della copia
(\texttt{quandoSalvato}: ``oggi alle 8:14'', ``ieri alle 21:03'', ``il 3
agosto''). Sono due i modi di trovarsi davanti alla copia, e il banner li
distingue: il dispositivo dice che non c'\`e connessione, oppure crede che ci
sia e il server non ha comunque risposto. Il secondo \`e quello che prima
sembrava un difetto dell'app.

Nella stessa direzione va una scelta sull'errore: con l'elenco salvato a
schermo, la riga rossa sparisce. Il banner sta gi\`a spiegando perch\'e niente
si muove, e un errore sotto direbbe la stessa cosa con una voce pi\`u
allarmante. Per la stessa ragione un fallimento di rete non genera pi\`u eventi
Sentry n\'e nel caricamento della lista n\'e in \texttt{assicuraAccount}: adesso
che l'app si apre in treno, quelle chiamate falliscono a ogni avvio, ed \`e
l'esito previsto, non un incidente.

\chapter{Promemoria locali}
\label{sec:notifiche}

La specifica escludeva le notifiche push, e restano escluse. Quelle
implementate sono \textbf{notifiche locali}: le programma il dispositivo per
un istante futuro e le mostra da s\'e, anche ad app chiusa. La differenza non
\`e terminologica ed \`e la ragione per cui la funzionalit\`a rientra nei
vincoli del progetto invece di violarli: non c'\`e nessun servizio push,
nessun token da registrare, nessun costo, e soprattutto nessun server che
debba sapere che cosa si ripassa e quando. Un promemoria push avrebbe
richiesto di mandare titolo e orario a un'infrastruttura terza; qui il titolo
non lascia mai il telefono.

\section{Come sono divise le responsabilit\`a}

\begin{center}
\small
\begin{tabular}{@{}p{3.5cm}p{4.3cm}p{6.7cm}@{}}
\toprule
\textbf{Livello} & \textbf{File} & \textbf{Che cosa decide} \\
\midrule
config &
  \texttt{notifications.ts} &
  L'impianto del dispositivo: permesso, canale Android (senza, su Android non
  si programma nulla), comportamento in primo piano. \\
\addlinespace
Model (puro) &
  \texttt{notificheLogic.ts} &
  Quali occorrenze meritino un promemoria, e che cosa sia cambiato dall'ultima
  volta. Nessun I/O, quindi la politica \`e testabile senza
  \texttt{expo-notifications}. \\
\addlinespace
Model (I/O) &
  \texttt{notificheRepo.ts} &
  Pianifica e cancella. Interfaccia pi\`u implementazione di default, come
  \texttt{RipassiRepo}. \\
\addlinespace
Controller &
  \texttt{useNotificheRipassi.ts} &
  Tiene allineato ci\`o che \`e programmato sul dispositivo con l'elenco che
  \texttt{RipassiProvider} ha gi\`a in mano. \\
\bottomrule
\end{tabular}
\end{center}

\section{Le decisioni che contano}

\begin{itemize}[nosep]
  \item \textbf{L'id della notifica \`e l'id dell'occorrenza}. Lo scheduler del
        sistema operativo \`e indicizzato come i dati dell'app, quindi
        ``riprogramma questa occorrenza'' \`e ``programma di nuovo con lo
        stesso id'' e non serve nessuna tabella di corrispondenza fra
        occorrenze e notifiche --- che sarebbe stato un terzo stato da tenere
        coerente con gli altri due.
  \item \textbf{Si programma solo il futuro non completato}. Un'occorrenza gi\`a
        passata non viene riprogrammata: il promemoria per ``adesso'' non
        significa niente quando ``adesso'' \`e trascorso, e un ripasso saltato
        si ritrova nello Storico, non in una coda di notifiche.
  \item \textbf{Un diff, non una riprogrammazione}. Un'occorrenza gi\`a
        programmata per lo stesso istante e con lo stesso titolo viene lasciata
        stare: riprogrammarla significherebbe cancellare e ricreare un
        promemoria gi\`a corretto a ogni ricaricamento innescato dal Realtime,
        cio\`e continuamente.
  \item \textbf{Il permesso si chiede quando c'\`e un motivo}, non all'avvio:
        finch\'e non c'\`e niente da ricordare non si chiede niente, come Drive
        chiede il consenso al primo allegato. E dopo un rifiuto non si richiede
        pi\`u --- quella decisione appartiene all'utente e alle impostazioni di
        sistema, non a una schermata che la ripropone a ogni ricarica.
  \item \textbf{Una data illeggibile viene esclusa}, non fatta cadere su
        ``adesso'' o su ``mai'': programmare un promemoria per un istante
        indecifrabile non \`e un ripiego sicuro in nessuna delle due direzioni.
  \item \textbf{In primo piano la notifica si vede lo stesso}. Niente
        nell'app dice gi\`a ``\`e ora'' a chi sta guardando la Home, quindi
        sopprimerla renderebbe il promemoria affidabile o no a seconda di un
        dettaglio che l'utente non controlla.
\end{itemize}

\chapter{Sicurezza}

\section{Autenticazione}
\begin{itemize}[nosep]
  \item \textbf{OAuth Google con PKCE} (non implicit flow): il redirect
        \texttt{ripassa://} trasporta solo un \texttt{?code=} opaco, scambiato
        per la sessione con un \emph{code verifier} noto solo al client. Un
        URL di redirect intercettato o loggato non \`e sufficiente a rubare
        la sessione.
  \item \textbf{Sessione cifrata a riposo}: JWT e refresh token in
        \texttt{expo-secure-store} (Keychain iOS / Keystore Android), con
        spezzettamento in voci da $\le$2~KB (limite di SecureStore). Prima
        erano in AsyncStorage, in chiaro sul filesystem.
  \item \textbf{Refresh in foreground}: listener AppState che riavvia
        l'auto-refresh del token quando l'app torna attiva (pattern
        raccomandato Supabase); elimina i re-login forzati.
  \item \textbf{La copia locale della sessione sta nello stesso posto
        dell'originale}: SecureStore, non un file in chiaro. Porta lo stesso
        refresh token, quindi vale la stessa protezione --- la ragione per cui
        esiste \`e nel cap.~\ref{sec:offline}, non qui, ma dove viene scritta
        \`e una decisione di sicurezza. Un valore troncato o estraneo viene
        rifiutato alla lettura: senza uno dei tre campi che la rendono
        utilizzabile (i due token e l'id utente) aprirebbe l'app in uno stato
        che niente pu\`o riparare. Viene cancellata al logout e quando il
        server dice che la sessione non c'\`e pi\`u, \textbf{mai} solo perch\'e
        il server non si \`e fatto sentire.
  \item Solo la chiave \emph{publishable} \`e nel client; la chiave
        \emph{secret} non compare da nessuna parte nel progetto.
\end{itemize}

\subsection*{Il redirect torna per due strade, e a volte a processo morto}

\`E la parte pi\`u difficile del sistema e merita un disegno. Il redirect che
chiude il login pu\`o raggiungere l'app in tre modi diversi, e per due di essi
il codice deve essere pronto a non essere l'unico ad averlo ricevuto.

\begin{figure}[H]
\centering
\resizebox{\textwidth}{!}{%
\begin{tikzpicture}[node distance=1.0cm]
  % attori
  \node[attore] (app) {App};
  \node[attore, right=1.6cm of app] (br) {Browser\\\scriptsize (custom tab)};
  \node[attore, right=1.6cm of br] (goog) {Google};
  \node[attore, right=1.6cm of goog] (sb) {Supabase};

  \foreach \n in {app,br,goog,sb} \draw[vita] (\n.south) -- ++(0,-7.9);

  \draw[msg] ($(app.south)+(0,-0.75)$) -- node[etichetta, above] {signInWithOAuth} ($(sb.south)+(0,-0.75)$);
  \draw[msgret] ($(sb.south)+(0,-1.35)$) -- node[etichetta, above] {url di consenso} ($(app.south)+(0,-1.35)$);
  \draw[msg] ($(app.south)+(0,-1.95)$) -- node[etichetta, above] {apre} ($(br.south)+(0,-1.95)$);
  \draw[msg] ($(br.south)+(0,-2.55)$) -- node[etichetta, above] {consenso} ($(goog.south)+(0,-2.55)$);
  \draw[msgret] ($(goog.south)+(0,-3.15)$) -- node[etichetta, above] {redirect \texttt{ripassa://?code=}} ($(br.south)+(0,-3.15)$);

  % Riquadro opaco: copre le linee di vita, che qui non aggiungono nulla.
  \node[nota, fill=white, draw=grigiobordo, rounded corners=2pt, inner sep=7pt,
        text width=12.4cm, anchor=north west]
       at ($(app.south)+(-1.3,-3.65)$) {%
    \textbf{Da qui il codice raggiunge l'app per una di tre strade:}\\[2pt]
    \textbf{(a)} risultato di \texttt{openAuthSessionAsync} --- il caso normale\\
    \textbf{(b)} evento \texttt{url} (deep link): Android consegna l'intent all'app ancora viva\\
    \textbf{(c)} \texttt{getInitialURL()}: il processo era stato ucciso, e il redirect lo riavvia};

  \draw[msg] ($(app.south)+(0,-6.35)$) -- node[etichetta, above] {exchangeCodeForSession(code)} ($(sb.south)+(0,-6.35)$);
  \draw[msgret] ($(sb.south)+(0,-7.1)$) -- node[etichetta, above] {sessione} ($(app.south)+(0,-7.1)$);

  \node[nota, text width=13.6cm, anchor=north west]
       at ($(app.south)+(-1.3,-8.15)$) {%
    Il codice \`e monouso, e (a) e (b) possono arrivare entrambi in ordine non
    garantito: \texttt{codiciUsati} riconosce il secondo e gli fa attendere
    l'esito del primo, invece di bruciare il codice una seconda volta.};
\end{tikzpicture}}
\caption{Login Google con PKCE. Il ramo (c) \`e quello che rendeva il difetto
invisibile: la promise che aspettava il browser moriva col processo, e l'app
riapriva la schermata di login come se non fosse mai successo niente.}
\end{figure}

L'autorizzazione a Drive percorre la stessa struttura con una differenza: il
\emph{code verifier} PKCE viveva solo dentro l'oggetto \texttt{AuthRequest}, e
quindi non sopravviveva alla morte del processo. Ora viene scritto in
SecureStore \emph{prima} di aprire il browser, insieme allo \texttt{state} che
fa da guardia CSRF sul ramo (c), dove non c'\`e pi\`u nessun
\texttt{AuthRequest} a controllarlo.

\section{Account e identit\`a di accesso}

L'isolamento \`e per \textbf{account}, non per riga di \texttt{auth.users}.
La distinzione non \`e formale: un utente di \texttt{auth.users} \`e
\emph{un modo di accedere}, non una persona. Registrarsi con email e password
e poi entrare con Google — stesso indirizzo — produce due righe distinte,
perch\'e l'identit\`a \`e la coppia \texttt{(provider, subject)} e l'email ne
\`e solo un attributo. Con la propriet\`a agganciata a \texttt{auth.users.id}
la stessa persona si ritrovava due insiemi di ripassi disgiunti.

Da qui due tabelle al posto di una: \texttt{account} (la persona, proprietaria
dei dati) e \texttt{identita} (un modo di entrare, N per account, uno per riga
di \texttt{auth.users}). Qualunque identit\`a risolve allo stesso account.

Un'identit\`a entra in un account esistente \textbf{solo a email verificata da
entrambe le parti}. Senza il vincolo il modello sarebbe peggiore del difetto
che risolve: chi registra in anticipo l'indirizzo di un altro con una password
scelta da s\'e si vedrebbe consegnare l'account al primo accesso Google della
vittima. Per questo le identit\`a \texttt{email} non si collegano mai alla
creazione, ma solo quando il link di conferma viene davvero cliccato.
L'invariante \`e imposta dal database — un indice unico parziale garantisce al
massimo un account \emph{verificato} per indirizzo — non dal codice
applicativo. Dettagli e casi limite in \texttt{docs/account-identita.md}.

\section{Row Level Security}

Tutte le tabelle hanno RLS attiva con isolamento per account. Su
\texttt{occorrenze} e \texttt{allegati} la policy verifica sia
\texttt{account\_id} sulla riga, sia — difesa in profondit\`a — che il ripasso
genitore appartenga allo stesso account:

\begin{lstlisting}[language=SQL]
create policy occorrenze_owner on public.occorrenze
  for all using (
    account_id = public.account_corrente()
    and exists (select 1 from public.ripassi r
                where r.id = ripasso_id
                  and r.account_id = public.account_corrente())
  ) with check ( -- identico, piu' user_id is null or = auth.uid() --);
\end{lstlisting}

\texttt{account\_corrente()} traduce la sessione nel suo account. \`E
\texttt{stable}, quindi Postgres la valuta una volta per statement e non una
per riga, e \texttt{security definer}, perch\'e altrimenti la RLS su
\texttt{identita} avrebbe bisogno di questa stessa funzione per decidere se la
riga \`e leggibile.

\section{I binari degli allegati}

Non sono su Supabase e quindi non sono soggetti a RLS: stanno su Google Drive.
Le garanzie sono di altra natura e vanno dette per quello che sono:

\begin{itemize}[nosep]
  \item \textbf{Scope minimo}: \texttt{drive.file}. L'app vede
        \emph{soltanto} i file che ha creato lei, non il resto del Drive
        dell'utente. Anche la lettura dell'account
        (\texttt{about.get}) resta dentro questo scope: chiedere
        \texttt{userinfo.email} solo per mostrare un indirizzo avrebbe
        allargato i permessi.
  \item \textbf{Token separati e cifrati a riposo}: access e refresh token
        Drive stanno in \texttt{expo-secure-store}, non vengono mai inviati a
        Supabase e vengono cancellati al logout.
  \item \textbf{Nessun URL pubblico}: con \texttt{drive.file} i file non sono
        raggiungibili da un link. Per mostrarli l'app li scarica in locale con
        il token, il che elimina del tutto la superficie ``URL firmato che
        circola''.
  \item \textbf{Confine di responsabilit\`a}: dei binari risponde Google, dei
        metadati Postgres con RLS. Il prezzo, dichiarato, \`e che l'utente pu\`o
        cancellare i file dal proprio Drive senza che l'app lo sappia --- la
        rotazione della cache lo scopre e ora lo segnala (cap.~\ref{sec:cache}).
\end{itemize}

\paragraph{Le due funzioni transazionali, e perch\'e non sono un buco.} Il
riordino degli allegati passa da \texttt{riordina\_allegati} e la
riprogrammazione delle date da \texttt{sposta\_occorrenze}. Entrambe sono
dichiarate in \texttt{SECURITY INVOKER} --- deliberatamente, ed \`e la sola
cosa che le rende innocue: girano con i privilegi di chi chiama, quindi restano
soggette alle policy \texttt{allegati\_owner} e \texttt{occorrenze\_owner}. Un
utente non pu\`o toccare righe che non gli appartengono nemmeno passando per la
funzione, ed \`e il motivo per cui nessuna delle due filtra per
\texttt{account\_id} nel proprio corpo: lo fa gi\`a la RLS, e ripeterlo
suggerirebbe che sia quel controllo a proteggere le righe. In
\texttt{SECURITY DEFINER} sarebbero invece due modi per aggirare l'isolamento
passando un elenco di id altrui.

\paragraph{Permessi richiesti sul dispositivo.} Solo la fotocamera
(\texttt{NSCameraUsageDescription}, \texttt{CAMERA}) e la galleria
(\texttt{NSPhotoLibraryUsageDescription}), cio\`e i due picker effettivamente
usati. Il permesso \texttt{RECORD\_AUDIO}, ereditato per copia da una
configurazione di esempio, \`e stato rimosso: l'app non ha nessuna funzione
audio, e chiedere il microfono a un'app di ripassi \`e esattamente il tipo di
richiesta che il principio del minimo privilegio esiste per evitare.

Alla revisione successiva \`e emerso che il paragrafo era incompleto:
\texttt{app.json} dichiarava anche \texttt{READ\_EXTERNAL\_STORAGE}, un
secondo residuo dello stesso tipo. Da Android~13 quel permesso \`e deprecato e
non d\`a pi\`u accesso ai media, e i picker usati qui non ne hanno comunque
bisogno --- \texttt{expo-image-picker} passa dal photo picker di sistema e
\texttt{expo-document-picker} da SAF, e entrambi restituiscono URI gi\`a
autorizzati. \`E stato rimosso: \texttt{app.json} oggi dichiara la sola
\texttt{CAMERA}.

A questi si \`e aggiunto il permesso di \textbf{mostrare notifiche}, che su
Android~13 e successivi \`e a richiesta esplicita. Non \`e dichiarato a mano:
lo porta \texttt{expo-notifications} nel proprio manifest, che l'unione dei
manifest fonde nell'app. Il criterio del minimo privilegio qui non si applica
alla dichiarazione ma al \emph{momento}: non si chiede all'avvio, ma la prima
volta che esiste davvero qualcosa da ricordare (cap.~\ref{sec:notifiche}), ed
\`e chiesto una volta sola.

\chapter{Dove vivono gli allegati}
\label{sec:drive}

\`E lo scostamento pi\`u vistoso rispetto alla specifica iniziale, ed \`e
deliberato. La sezione~2 della specifica poneva fra le decisioni vincolanti
``nessuna dipendenza da iCloud o Google Drive per lo storage applicativo'';
oggi i binari degli allegati stanno proprio sul Google Drive dell'utente.
Questo capitolo esiste perch\'e una decisione vincolante ribaltata senza
spiegazione \`e indistinguibile da un requisito disatteso.

\section{Il vincolo che ha forzato la scelta}

Il piano gratuito Supabase d\`a 1~GB di file storage. La specifica lo aveva
gi\`a individuato come ``l'unico vincolo che collide con \emph{nessun limite di
dimensione}'' e stimava il superamento in qualche mese di uso reale con foto e
PDF di studio. Le alternative valutate:

\begin{center}
\small
\begin{tabular}{@{}p{4.3cm}p{4.0cm}p{6.2cm}@{}}
\toprule
\textbf{Opzione} & \textbf{Costo} & \textbf{Perch\'e s\`i / no} \\
\midrule
Restare su Supabase Storage & gratis fino a 1~GB & Il tetto arriva, e arriva
  presto. Superarlo significa perdere l'accesso ai propri appunti. \\
\addlinespace
Supabase Pro & 25~\$/mese & Fuori dal budget dichiarato (\textasciitilde2~\texteuro/mese)
  di un ordine di grandezza. \\
\addlinespace
Compressione pi\`u aggressiva & gratis & Sposta il problema di qualche mese e
  peggiora la leggibilit\`a di appunti fotografati, che \`e il contenuto
  principale. \\
\addlinespace
\textbf{Google Drive dell'utente} & \textbf{gratis (15~GB)} & \textbf{Scelta.}
  Quindici volte lo spazio, sullo stesso account Google gi\`a usato per il
  login, senza nuovi servizi da gestire. \\
\bottomrule
\end{tabular}
\end{center}

\section{Perch\'e non riporta il problema di partenza}

L'obiezione ovvia: l'app nasce per risolvere la frammentazione di ``Genio in
21 Giorni'', che su iPad usava iCloud e su Android Google Drive. La
distinzione \`e che il difetto non era \emph{usare Drive}, era usare
\textbf{archivi diversi su piattaforme diverse}. Un unico Drive, raggiunto con
le stesse credenziali da Android, iPad e Mac, non frammenta nulla: \`e un solo
archivio con un solo account, esattamente ci\`o che la decisione vincolante
voleva ottenere. Restano su Supabase tutti i metadati --- titoli, note, date,
ordine, nomi --- quindi la sincronizzazione in tempo reale e la logica di
dominio non cambiano di una riga.

\section{Cosa cambia, in concreto}

\begin{itemize}[nosep]
  \item \textbf{Un secondo consenso OAuth}, distinto dal login. L'utente lo
        concede la prima volta che allega un file, non all'avvio. \`E il costo
        pi\`u alto della scelta: un secondo flusso di autorizzazione con i
        suoi modi di fallire, che si porta dietro \texttt{driveAuth.ts},
        \texttt{useDriveAuth.ts} e due regole dedicate nella traduzione degli
        errori.
  \item \textbf{Lo spazio \`e quello dell'utente}, non del progetto: non c'\`e
        pi\`u un tetto condiviso da sorvegliare. La compressione resta comunque
        attiva, per rispetto dei 15~GB altrui.
  \item \textbf{L'utente pu\`o cancellare i file dal proprio Drive} senza che
        l'app lo sappia. La rotazione della cache se ne accorge e ora lo
        segnala invece di ignorarlo (cap.~\ref{sec:cache}).
  \item \textbf{Serve una build installata}: il redirect OAuth di Google per i
        client nativi dev'essere \texttt{<applicationId>:/oauthredirect}, che
        Expo Go non pu\`o registrare.
  \item \textbf{Il campo si chiama ancora \texttt{storage\_path}}. Contiene un
        ID Drive, non un percorso. Il nome \`e rimasto per non imporre una
        migrazione dello schema a un dato che \`e comunque opaco; il tipo lo
        documenta esplicitamente.
\end{itemize}

\chapter{Ciclo di vita di un allegato}

\section{Caricamento}

Due scritture su due sistemi che non condividono nessuna transazione, quindi
la parte interessante \`e la compensazione.

\begin{enumerate}[nosep]
  \item Se \`e un'immagine viene compressa (1600px sul lato lungo, JPEG 70\%).
        Un PDF passa invariato.
  \item \textbf{Upload su Drive in due passi}: prima i metadati (crea un file
        vuoto nella cartella \texttt{ripassiProgrammati} e restituisce l'ID),
        poi il contenuto binario via \texttt{FileSystem.uploadAsync}. Due passi
        e non uno per non caricare l'intero file in memoria come base64: una
        foto compressa resta comunque di centinaia di KB, e la conversione era
        pura spesa. Se il secondo passo fallisce, il file vuoto appena creato
        viene cancellato.
  \item \textbf{Riga in \texttt{allegati}} con l'ID Drive in
        \texttt{storage\_path}. Se l'insert fallisce, il file su Drive viene
        cancellato: un binario senza metadati sarebbe spazzatura invisibile
        nel Drive dell'utente.
  \item \texttt{size\_bytes} registra la dimensione del file \emph{caricato},
        non quella dichiarata dal picker. Erano due numeri diversi: le immagini
        vengono compresse, e salvare l'originale sovrastimava il consumo di
        circa un ordine di grandezza --- proprio sulla classe di file che
        domina il volume, e proprio mentre il consumo di spazio \`e la cosa da
        tenere d'occhio.
\end{enumerate}

Il caricamento \textbf{non viene mai ritentato}: crea un file nuovo e una riga
nuova, quindi non \`e idempotente (cap.~\ref{sec:retry}).

\section{Apertura}

Un allegato si apre sempre da un file \emph{locale}: con lo scope
\texttt{drive.file} non esiste un URL pubblico da passare a un viewer.

\begin{enumerate}[nosep]
  \item Se \`e nella cache della finestra, si usa quella copia (dopo aver
        verificato che il file esista ancora davvero sul dispositivo).
  \item Altrimenti si scarica in un file temporaneo nella cache di sistema.
        Questo download \`e idempotente, quindi \emph{viene} ritentato.
  \item Le immagini si mostrano in un viewer in-app a schermo intero; PDF e
        altri formati vanno al viewer di sistema (intent \texttt{VIEW} con
        \texttt{content://} su Android, share sheet / Quick Look su iOS).
\end{enumerate}

\emph{Mai} un \texttt{file://} verso WebBrowser: su Android \`e una
\texttt{FileUriExposedException}.

\paragraph{Ciclo di vita dei temporanei.} Il punto~2 aveva un difetto che la
prima stesura di questo capitolo non nominava: quei file non venivano
rimossi da niente. Il nome deriva dall'id dell'allegato, quindi riaprire lo
\emph{stesso} allegato sovrascrive, ma ogni allegato \emph{distinto} aperto
fuori finestra --- cio\`e la normale consultazione dello storico --- lasciava
dietro di s\'e un PDF o una foto per sempre. La cartella di cache di sistema
viene potata dal sistema operativo solo sotto pressione di spazio, e su
Android quel momento tende ad arrivare ben dopo che l'utente si \`e accorto
dell'app che occupa centinaia di MB. In un progetto che ha eletto il consumo
di spazio a vincolo dimensionante (cap.~\ref{sec:drive}) era una falla di
categoria.

Ora i temporanei stanno in una sottocartella dedicata
(\texttt{allegati-tmp/}), il che rende la potatura l'elenco di una cartella
che non contiene altro invece di un filtro su una convenzione di nomi. La
politica \`e a tempo: si conservano sette giorni dall'ultimo uso, la
potatura gira insieme alla rotazione della cache (al massimo una volta al
giorno, cadenza adatta a una politica misurata in giorni) e al logout si
cancella tutto, perch\'e a quel punto quei file non riguardano pi\`u nessuno.
La selezione \`e una funzione pura in \texttt{cacheLogic.ts}, testata: un file
di cui non si riesce a leggere la data conta come scaduto, perch\'e non
poterlo dimostrare recente e tenerlo per sempre \`e esattamente l'esito da
evitare. Il costo di un errore di stima \`e un download in pi\`u.

\section{Rinomina, riordino, eliminazione}

Rinomina e riordino toccano solo i metadati. Il riordino \`e una singola
chiamata transazionale (\texttt{riordina\_allegati}): la versione precedente
mandava un UPDATE per riga in parallelo, e un fallimento a met\`a lasciava
\texttt{order\_index} duplicati o con buchi, senza modo di tornare indietro.
L'eliminazione rimuove la riga, poi il file su Drive, poi la copia in cache --- in
quest'ordine: cancellare prima il binario lascerebbe una riga che punta a un
file inesistente, cio\`e un allegato che non si apre pi\`u.

La conseguenza dell'ordine giusto \`e per\`o che il fallimento del secondo
passo produce un orfano \emph{irrecuperabile}: la riga che conteneva l'id
Drive non esiste pi\`u, quindi n\'e l'app n\'e l'utente potranno mai risalire
a quale file nella cartella sia spazzatura, e lo spazio resta occupato per
sempre. Il codice lo ignorava in silenzio, con un \texttt{catch} vuoto, nello
stesso file che venti righe pi\`u su tratta l'orfano opposto (binario senza
metadati) come inaccettabile e implementa la compensazione. Ora
\texttt{elimina} restituisce un esito: l'operazione riesce lo stesso, perch\'e
per l'utente l'allegato \`e sparito davvero e un avviso sarebbe rumore, ma
l'esito riporta l'id rimasto orfano e il Controller genera l'evento Sentry
corrispondente. Stesso trattamento per il rollback fallito del caricamento,
che lasciava un orfano speculare.

\chapter{Affidabilit\`a in rete}

Tre problemi emersi rileggendo il codice in vista di una distribuzione ad
utenti reali, e come sono stati chiusi.

\section{Errori comprensibili}
\label{sec:errori}

Gli errori di Supabase/Postgres sono in inglese e citano dettagli interni
(\texttt{duplicate key value violates unique constraint "ripassi\_pkey"}).
Mostrarli in un \texttt{Alert} a un utente italiano \`e inaccettabile.
\texttt{errorMessages.ts} li classifica per categoria (rete, autenticazione,
permessi, duplicato, non trovato, spazio) e restituisce titolo e messaggio in
italiano, con un flag \texttt{ritentabile}.

La mappatura \`e una \textbf{tabella di regole}, non una catena di \texttt{if}:
le regole sono dati, il confronto \`e un solo ciclo, e aggiungere un caso
significa aggiungere una riga invece di trovare il punto giusto in una
funzione di 150 righe.

\paragraph{La proprietà critica \`e l'ordine.} Le regole si applicano dall'alto
in basso, quindi ci\`o che tiene corretto il modulo non \`e ``esiste la regola
giusta'' ma ``nessuna regola pi\`u larga arriva prima''. Due difetti reali
nascevano proprio l\`i, e sono stati chiusi rendendo strutturale ci\`o che era
testuale:

\begin{itemize}[nosep]
  \item \textbf{Codici di stato}. Erano cercati come sottostringhe
        (\texttt{"500"}, \texttt{"404"}, \texttt{"413"}) dentro un testo che
        include il corpo della risposta di Google. Un errore \texttt{403} il cui
        JSON citava una quota di ``500 richieste'' veniva classificato guasto
        transitorio del server, e quindi \emph{ritentato} tre volte con
        backoff contro una richiesta che non poteva riuscire. Ora gli errori
        Drive trasportano lo stato come campo (\texttt{DriveHttpError.status})
        e le regole confrontano un numero.
  \item \textbf{Il frammento \texttt{"session"}} catturava qualunque messaggio
        di \texttt{expo-auth-session}: uno schema di redirect non registrato
        diceva all'utente ``la sessione \`e scaduta, esci e accedi di nuovo'',
        cio\`e l'unica azione che non poteva servire a niente. Ora i frammenti
        sono espliciti (\texttt{session expired}, \texttt{auth session
        missing}, \dots).
\end{itemize}

I test verificano la categoria, che il testo tecnico \emph{non} compaia nel
messaggio finale, e --- con un blocco dedicato di casi negativi --- che
nessuna regola larga preceda quelle specifiche.

\paragraph{L'unica eccezione, e perch\'e esiste.} Il messaggio tradotto \`e
sempre la voce principale, ma quando la traduzione \emph{non} riesce a
classificare l'errore (categoria \texttt{sconosciuto}) il testo originale viene
allegato, troncato, sotto la voce ``Dettagli''. Il motivo \`e pratico: un
caricamento di allegato pu\`o fallire su Drive, su Postgres o sul filesystem,
e ``Operazione non riuscita'' li rende indistinguibili --- l'utente non ha
niente da riferire e chi legge il crash report nemmeno. Vale in tre punti:
\texttt{mostraErrore} per la categoria sconosciuta, e i due messaggi del login
Google, che indicano anche in quale met\`a del flusso si \`e rotto e quale
redirect era atteso. Un errore \emph{riconosciuto} resta invece pulito: il
dettaglio tecnico non lo tocca mai.

\section{Offline esplicito}

Senza rilevamento della connettivit\`a, ``sei offline'' e ``\`e successo un
errore'' finivano nello stesso alert generico. \texttt{useConnettivita}
incapsula NetInfo dietro il Controller e restituisce un solo booleano; Home e
Login mostrano un banner dedicato. Solo una lettura esplicitamente negativa
conta come offline, cos\`i una sonda di raggiungibilit\`a lenta non produce un
falso allarme.

\section{Retry: dove s\`i e dove no}
\label{sec:retry}

\texttt{retry.ts} ritenta con backoff esponenziale, ma \textbf{solo} quando
l'errore \`e transitorio (rete, 5xx) — ritentare un errore di vincolo o di
permesso ripete solo lo stesso fallimento — e \textbf{solo} su operazioni
idempotenti.

\begin{center}
\small
\begin{tabular}{@{}p{5.0cm}p{1.3cm}p{8.2cm}@{}}
\toprule
\textbf{Operazione} & \textbf{Retry} & \textbf{Perch\'e} \\
\midrule
lettura elenco & s\`i & idempotente \\
modifica / completa & s\`i & stesso input, stesso stato finale \\
sposta occorrenze & s\`i & le date scritte sono assolute, non relative \\
elimina (ripasso, allegato) & s\`i & idempotente \\
rinomina / riordina allegati & s\`i & idempotente (il riordino \`e transazionale) \\
download per l'apertura & s\`i & idempotente, e l'utente sta aspettando \\
\midrule
\textbf{crea ripasso} & \textbf{no} & insert non idempotente \\
\textbf{carica allegato} & \textbf{no} & crea file Drive + riga nuovi \\
\textbf{download in rotazione cache} & \textbf{no} & scelta diversa: vedi sotto \\
\bottomrule
\end{tabular}
\end{center}

Le esclusioni sono deliberate e sono il punto meno ovvio del capitolo.

\textbf{Le prime due} per non idempotenza: un errore di rete pu\`o significare
``la richiesta \`e arrivata ma la risposta si \`e persa'', e ritentare
creerebbe un secondo ripasso, o un secondo file su Drive con una seconda riga
in \texttt{allegati} (il nome su Drive \`e generato con \texttt{Date.now()},
quindi ogni tentativo produce un file diverso). Chiedere all'utente di
ripremere \`e meno grave di un duplicato silenzioso.

\textbf{La terza per un motivo opposto}: il download \emph{sarebbe}
idempotente, ed \`e infatti ritentato quando l'utente apre un allegato. Ma la
rotazione della cache pu\`o riguardare decine di file in una volta, all'avvio,
tipicamente proprio quando la rete \`e cattiva: tre tentativi con backoff su
ciascuno bloccherebbero l'app per minuti nel momento peggiore. La rotazione
riprova alla successiva apertura, che \`e la stessa cosa senza l'attesa. Ci\`o
che invece non deve pi\`u succedere \`e che il fallimento passi inosservato:
vedi il capitolo sulla cache.

\subsection*{Un ritento invisibile \`e indistinguibile da un blocco}

\texttt{conRetry} aspetta fra un tentativo e l'altro con attese che
raddoppiano: fino a tre tentativi sono diversi secondi in cui a schermo non
succede niente. Dal lato di chi guarda \`e un'app bloccata, e la reazione
ragionevole davanti a un'app bloccata \`e toccare di nuovo --- cio\`e l'unica
cosa che non aiuta, e che su un'operazione non idempotente sarebbe pure peggio
di non aiutare.

La logica di ritento non cambia di una riga: resta nel Model, pura e senza
stato. Bastava usare il gancio \texttt{onRitento} che gi\`a esportava, ed \`e
quello che fa \texttt{useRitento}, un hook di Controller condiviso da
\texttt{useRipassi} e \texttt{useAllegati} --- gli unici due punti da cui
\texttt{conRetry} viene chiamato. Due dettagli non ovvi:

\begin{itemize}[nosep]
  \item \textbf{Un contatore, non un booleano}: pi\`u operazioni possono essere
        in volo insieme (una scrittura e il ricaricamento che ne segue), e con
        un booleano la prima che finisce spegnerebbe l'indicatore mentre la
        seconda sta ancora aspettando.
  \item \textbf{Si conta l'operazione, non il tentativo}: \texttt{onRitento}
        scatta a ogni tentativo, e contarli tutti lascerebbe il contatore sopra
        lo zero per sempre.
\end{itemize}

Lo stato esposto si chiama \texttt{ritentando} ed \`e distinto sia da
\texttt{loading} sia da \texttt{busy}: non \`e ``la lista non \`e mai
arrivata'' n\'e ``sto caricando'', \`e ``\`e gi\`a andata storta una volta e ci
sto riprovando''. Le tre schermate che avevano gi\`a un indicatore visibile gli
affiancano una riga.

\section{Ordinamento totale}

La classificazione attivi/storico viveva in \texttt{HomeScreen} e ordinava con
un fallback \texttt{?? ""} che, su una data assente, produce
\texttt{new Date("").getTime()} = \texttt{NaN} e rende il comparatore
inconsistente. Nel flusso attuale quel ramo non era raggiungibile (il filtro
\texttt{isStorico} garantisce che gli attivi abbiano sempre un'occorrenza
futura), ma la correttezza dipendeva da un invariante non locale, cio\`e dalla
coerenza fra due funzioni separate. La logica \`e stata spostata in
\texttt{ripassiLogic.ts} con una chiave di ordinamento totale (mai
\texttt{NaN}, con fallback su \texttt{created\_at}) e un tie-break sull'id,
cos\`i due ripassi con la stessa data non si scambiano di posto fra un
ricaricamento e l'altro.

\chapter{Crash reporting e resilienza}

Prima di distribuire build a utenti reali servivano due cose che l'MVP non
aveva: sapere \emph{che} un crash \`e avvenuto, ed evitare che un errore di
render lasci lo schermo bianco.

\begin{itemize}[nosep]
  \item \textbf{Sentry} (\texttt{@sentry/react-native}), incapsulato in
        \texttt{src/config/crashReporting.ts}. Vale la stessa regola del
        client Supabase: \emph{nessun altro modulo importa l'SDK}, si passa
        da \texttt{initCrashReporting()} e \texttt{reportError()}. Il DSN si
        legge con la stessa priorit\`a env $\rightarrow$
        \texttt{app.json/extra} usata altrove
        (\texttt{EXPO\_PUBLIC\_SENTRY\_DSN}).
  \item \textbf{Degrado morbido}: senza DSN il modulo fa no-op e logga un
        warning; l'app parte e builda normalmente. Il reporting \`e inoltre
        disattivato in \texttt{\_\_DEV\_\_}, per non mescolare errori di
        sviluppo e crash di produzione. \texttt{tracesSampleRate} \`e 0: qui
        interessano i crash, non il performance tracing.
  \item \textbf{ErrorBoundary} (\texttt{src/view/components/ErrorBoundary.tsx}),
        avvolge l'intero albero in \texttt{App.tsx}. Cattura gli errori di
        render/lifecycle, li inoltra a Sentry con lo stack dei componenti e
        mostra un fallback a tema con un pulsante ``Riprova'' che resetta lo
        stato, invece di un crash silenzioso. Deve essere una classe: React
        non espone gli error boundary come hook.
  \item \textbf{\texttt{Sentry.wrap}} sull'export di \texttt{App.tsx}, per
        intercettare anche ci\`o che sta fuori dal render tree (crash nativi,
        errori JS non gestiti).
\end{itemize}

Il config plugin \texttt{@sentry/react-native} \`e registrato in
\texttt{app.json}: le build native lo includono senza passaggi manuali. Resta
da registrare \texttt{EXPO\_PUBLIC\_SENTRY\_DSN} su EAS prima della
distribuzione (procedura in \texttt{BUILD.md}).

\section{La met\`a che mancava: un errore gestito raccontato da chi l'ha visto}

Tutto ci\`o che sta sopra riguarda i crash \emph{non} gestiti: ErrorBoundary e
il wrap sulla radice. Restava fuori esattamente la classe di guasti in cui si
inciampa davvero --- un upload su Drive che non \`e arrivato, una scrittura
rifiutata dal server: l'app li cattura, li traduce, li mostra, e l\`i
finiscono. L'utente li legge, nessun altro ne sa niente.

Il pulsante \textbf{``Segnala problema''} nel Profilo \`e la met\`a mancante.
Chi ha visto il problema scrive che cosa \`e successo, e la segnalazione porta
con s\'e il contesto che altrimenti dovrebbe ricostruire a memoria: l'ultimo
errore che l'app gli ha mostrato, l'indirizzo dell'account, piattaforma e
versione. Mezz'ora dopo ci si ricorda ``non si caricava'', non quale operazione
sia fallita n\'e che cosa avesse risposto Drive.

\begin{itemize}[nosep]
  \item \textbf{L'errore \`e mostrato nel form, non allegato di nascosto}: ci\`o
        che parte dal dispositivo \`e roba che l'utente ha letto.
  \item \textbf{Chi tiene l'ultimo errore \`e \texttt{avvisoErrore.ts}}, che \`e
        gi\`a l'unica strada per cui un errore arriva a schermo. \`E una
        variabile di modulo e non stato React: non c'\`e niente da
        ri-renderizzare quando cambia, e un Context avrebbe significato un
        provider attorno a tutta l'app per un valore che legge solo questo
        form, una volta, quando lo si apre.
  \item \textbf{L'invio sta in \texttt{crashReporting.ts}}, che resta l'unico
        file a importare l'SDK: vale qui la stessa regola del client Supabase.
  \item \textbf{Tre esiti, non un booleano}: ``le segnalazioni non sono attive
        in questa build'' e ``non \`e partita'' sono lo stesso fallimento per il
        codice e due frasi diverse per l'utente, che sulla seconda pu\`o agire e
        sulla prima no.
  \item \textbf{La conferma passa da un \texttt{flush}}: \texttt{captureMessage}
        si limita a mettere l'evento in coda, e senza aspettare che la coda si
        svuoti la schermata direbbe ``inviata'' a chi \`e in galleria.
\end{itemize}

Senza DSN --- cio\`e in ogni build di sviluppo --- resta un no-op con un
warning, come gi\`a fa \texttt{initCrashReporting}.

\section{Le domande a cui l'interfaccia da sola non sa rispondere}

Nel Profilo, sopra il pulsante di segnalazione, sta un \textbf{Aiuto} a domande
e risposte. La posizione \`e l'argomento: buona parte di ci\`o che si
segnalerebbe ha una risposta l\`i sopra.

Le domande coperte sono quelle che l'app non riesce a spiegare mentre fa le
cose, e sono anche i punti su cui una segnalazione arriverebbe scritta come
``non funziona'': perch\'e Drive chiede il consenso e poi lo richiede, perch\'e
i ripassi sono gli stessi entrando con la password o con Google, che cosa resta
leggibile senza connessione, che cosa succede a un allegato che non \`e
arrivato, che cosa comporta uscire dall'account. \`E testo statico, nessuna
chiamata: un elenco a tendina, una domanda aperta alla volta --- con tutte
aperte la sezione diventa un muro di testo in cui la domanda che interessa non
si trova pi\`u. Il test non ricopia le risposte (fallirebbe a ogni virgola
corretta) ma tocca ciascuna delle domande che sono il motivo della sezione: se
una sparisce dall'elenco, se ne accorge.

La riga da premere con il suo chevron \`e diventata un componente condiviso
(\texttt{TestataATendina}): la striscia ``Come funziona?'' della Home e ogni
domanda dell'Aiuto sono lo stesso gesto con la stessa affordance, e stavano per
essere disegnate due volte. Il chevron \textbf{ruota} invece di cambiare
simbolo: in molti font il carattere ``su'' e quello ``gi\`u'' hanno larghezze
diverse, e il titolo accanto si sposterebbe a ogni apertura.

\chapter{Test e qualit\`a}
\label{sec:test}

\begin{lstlisting}
npm run verifica   # lint + typecheck + test, in quest'ordine
npm run lint       # ESLint: regole Expo + confini fra i livelli
npm run typecheck  # tsc --noEmit (strict)
npm test           # 485 test su 38 suite (jest-expo)
npm run test:coverage
\end{lstlisting}

Gli stessi comandi girano in integrazione continua
(\texttt{.github/workflows/ci.yml}) a ogni push e su ogni pull request. Nessuno
tocca la rete: la suite copre logica pura e moduli il cui I/O \`e simulato,
quindi il job \`e autosufficiente e non richiede segreti. La CI esegue la
suite \emph{una volta sola}, con la copertura: prima la eseguiva due volte,
una in pi\`u che non aggiungeva nessuna informazione.

\section{La seconda suite: i test SQL}

Esiste una seconda suite che il capitolo precedente non nominava, e che vale
la pena nominare perch\'e copre la parte del sistema con le conseguenze pi\`u
severe: \texttt{supabase/tests/} verifica su righe vere il modello
account/identit\`a e le policy RLS --- cio\`e il meccanismo che impedisce a un
account di leggere i ripassi di un altro, e il vincolo di email verificata che
impedisce a un indirizzo occupato di catturare l'account della vittima. Sono
propriet\`a che non si vedono leggendo una definizione di tabella: vanno
provocate.

\begin{lstlisting}
npm run test:db    # richiede SUPABASE_DB_URL
\end{lstlisting}

Non girano in integrazione continua, e la ragione \`e dichiarata: sono scritti
contro lo schema \texttt{auth} di un progetto Supabase reale, non contro un
Postgres nudo, e terminano in \texttt{rollback} in modo da non lasciare
traccia. Portarli in CI richiederebbe lo stack Supabase locale completo. Vanno
eseguiti a mano dopo ogni modifica allo schema o alle policy: un test che non
gira \`e documentazione, non verifica, e questo \`e il limite pi\`u serio che
resta nella dimensione dei test.

\section{Copertura per livello}

\begin{center}
\small
\begin{tabular}{@{}lrp{8.4cm}@{}}
\toprule
\textbf{Livello} & \textbf{Soglia} & \textbf{Che cosa \`e coperto} \\
\midrule
\texttt{model} & 85\,\% & logica di dominio, repository (Supabase e Drive
  simulati), cache SQLite, istantanea offline e sessione locale, traduzione
  errori, retry, token Drive, logica dei promemoria \\
\texttt{config} & 62\,\% & risoluzione della configurazione, DSN Sentry,
  client ID Drive, guardia di init, permesso e canale delle notifiche \\
\texttt{controller} & 55\,\% & gli hook principali; \texttt{useAuth} resta
  coperto solo in parte (vedi sotto) \\
\texttt{view} & 79\,\% & helper puri (date, calendario) e tutte e cinque le
  schermate \\
\midrule
\textbf{globale} & \textbf{72\,\%} & rete di sicurezza sotto le quattro
  soglie di livello \\
\bottomrule
\end{tabular}
\end{center}

La tabella riporta le \emph{soglie} di \texttt{package.json} e non le
percentuali misurate. La differenza non \`e formale: una soglia \`e applicata
dalla CI e non pu\`o divergere dal codice, una percentuale trascritta a mano
invecchia al primo commit --- come era gi\`a successo al conteggio dei test,
che questa relazione riportava con tre valori diversi.

Le soglie sono per livello, e non una sola media globale, perch\'e i quattro
livelli hanno rischi diversi. Con la sola soglia globale la copertura alta
del \texttt{model} faceva da cuscinetto a quella bassa della \texttt{view}: si
potevano cancellare i test di un hook e restare sopra soglia aggiungendo
qualche test di componente. Una media non \`e una garanzia.

\paragraph{Una soglia troppo bassa non \`e prudenza.} La \texttt{view} \`e
passata da 29\,\% a 79\,\% e la globale da 60\,\% a 72\,\%, ed \`e questo il
punto della revisione, non un effetto collaterale: quando la copertura
realmente misurata era all'84\,\%, una soglia al 29\,\% non proteggeva da
nessuna regressione che non fosse la cancellazione di met\`a suite. I valori
restano deliberatamente qualche punto sotto il misurato --- bloccano una
regressione senza far fallire la CI al primo file nuovo --- ma la distanza
dev'essere di qualche punto, non di cinquanta.

\section{Che cosa i test verificano davvero}

Le asserzioni sono sul comportamento, non sull'implementazione. Alcuni esempi
che valgono pi\`u dell'elenco completo:

\begin{itemize}[nosep]
  \item \textbf{Regressione da perdita di dati}: aprendo la scheda di un
        ripasso prima che la lista sia caricata, un tocco su Salva non deve
        scrivere campi vuoti sopra titolo e note memorizzati.
  \item \textbf{Non idempotenza}: un errore di rete durante la creazione di un
        ripasso o il caricamento di un allegato deve produrre \emph{una sola}
        chiamata, mai due.
  \item \textbf{Idempotenza}: rinomina, riordino ed eliminazione devono invece
        ritentare, e ricaricare la lista una volta sola.
  \item \textbf{Compensazione}: se l'insert dei metadati fallisce, il file
        appena creato su Drive dev'essere cancellato; e un rollback fallito
        non deve mascherare l'errore originale.
  \item \textbf{Ordine delle regole di traduzione}: casi negativi che fissano i
        confini fra regole (un 403 con ``500'' nel corpo non \`e un guasto
        transitorio; un errore di \texttt{expo-auth-session} non \`e una
        sessione scaduta).
  \item \textbf{Osservabilit\`a della cache}: un download fallito non ferma gli
        altri, viene contato, e genera un evento Sentry solo se non \`e
        assenza di rete.
  \item \textbf{Ordinamento dei figli}: Postgres restituisce le righe innestate
        in ordine arbitrario; il repository garantisce occorrenze cronologiche
        e allegati nell'ordine scelto dall'utente.
  \item \textbf{Confini delle schermate}, non percorsi felici: che ``Elimina''
        ed ``Esci'' passino da una conferma; che il form resti aperto quando il
        Controller dice di non aver finito (il ripasso \`e salvato ma un
        allegato no, e chiudere farebbe sparire i file scelti); che una freccia
        di riordino non scriva niente sui due estremi; che un nome vuoto non
        sovrascriva quello che c'era; che un'apertura o uno spostamento falliti
        non restino muti --- il modale si chiude comunque, e senza avviso non
        ci sarebbe modo di distinguere ``non ha funzionato'' da ``non ho
        premuto bene''.
  \item \textbf{Copertura della tabella delle regole d'errore}: il test parte
        dalla tabella e chiede a ciascuna voce di essere raggiungibile da
        almeno un campione, confrontando l'\emph{identit\`a} dell'oggetto
        restituito e non il titolo (due regole possono produrre lo stesso
        titolo, e confrontare i testi direbbe ``coperta'' per una regola in
        realt\`a mai raggiunta perch\'e un'altra la precede). Il controllo
        simmetrico --- nessun campione inerte --- impedisce alla lista di
        campioni di sopravvivere alla regola che copriva. Una regola aggiunta
        in mezzo senza un caso che la raggiunga fa fallire la suite, che \`e
        esattamente il cambiamento che nessuno noterebbe.
  \item \textbf{Apertura offline}: che una sessione salvata decida il primo
        frame, che un \texttt{null} dovuto alla rete non venga scambiato per un
        logout, e che un'istantanea letta dal disco non sovrascriva righe che
        il server ha gi\`a confermato.
  \item \textbf{Promemoria}: che il canale Android venga creato (senza, non si
        programma nulla) e che su iOS non ci si provi; che un rifiuto del
        permesso non venga richiesto di nuovo a ogni ricarica; che il
        promemoria vada in coda sotto l'id dell'occorrenza stessa, che \`e
        ci\`o che rende superflua una tabella di corrispondenza.
\end{itemize}

\section{Limiti noti di copertura}

Dichiarati, non nascosti:

\begin{itemize}[nosep]
  \item \texttt{useAuth} e \texttt{useDriveAuth} restano coperti solo in parte:
        \`e il \emph{viaggio attraverso il browser} a non essere simulato. Ci\`o
        che \`e deterministico ha ora una suite --- per \texttt{useAuth}, il
        collegamento di Google (codice scambiato, fallimento non scambiato per
        riuscita perch\'e una sessione esisteva gi\`a, browser chiuso
        dall'utente, rifiuto tradotto) e tutta l'apertura offline descritta al
        cap.~\ref{sec:offline}. Sono il codice pi\`u delicato del progetto, ma
        anche quello il cui comportamento dipende da come il sistema operativo
        consegna un redirect e da quando decide di uccidere il processo:
        simularlo darebbe fiducia in un modello, non nel dispositivo. Restano
        verificati a mano, sui casi documentati in \texttt{BUILD.md}.

        La formulazione precedente si fermava qui, e diceva troppo: era vera
        per \emph{quando} il redirect arriva e se il processo sopravvive, non
        per che cosa il codice fa con lo \texttt{state} che riceve. La
        distinzione ora \`e esplicita, e la seconda met\`a \`e coperta:
        \texttt{driveAuth.ts} ha una suite propria su ci\`o che \`e
        deterministico --- persistenza del verifier PKCE prima di aprire il
        browser, guardia CSRF sullo \texttt{state}, redirect URI derivato
        dall'\texttt{applicationId}, unicit\`a del codice, e i tre esiti del
        refresh (valido, revocato, rete caduta). Resta fuori solo il viaggio
        attraverso il browser.

        \`E stato lo scrivere quei test a far emergere un difetto reale: la
        guardia CSRF richiedeva che lo \texttt{state} fosse presente
        \emph{sia} nel redirect \emph{sia} nel record salvato, quindi un
        redirect che semplicemente non ne portava nessuno saltava il
        controllo --- cio\`e esattamente ci\`o che farebbe chi sceglie che
        cosa mandare. Ora la condizione guarda solo lo \texttt{state} salvato.
  \item Il fallback della configurazione su \texttt{app.json/extra} non \`e
        coperto: jest-expo sostituisce \texttt{expo-constants} a livello di
        modulo nativo, quindi sotto jest \texttt{Constants.expoConfig} \`e
        \texttt{undefined} e non \`e sovrascrivibile con \texttt{jest.mock}. Il
        ramo \`e esercitato nelle build reali; il limite \`e annotato anche nel
        file di test.
  \item \emph{Chiuso in questa revisione.} La formulazione precedente diceva
        che le schermate Form e Dettaglio allegati non avevano test di render,
        perch\'e ci\`o che restava era ``disposizione grafica''. Non era vero
        --- e le quattro schermate a copertura zero (Form, Allegati, Login,
        Profilo) erano anche quelle che si toccano ogni volta che si aggiunge
        una funzione, quindi un refactor l\`i dentro si faceva al buio. Ora
        hanno una suite ciascuna, nello stile di quella della Home: i
        Controller sostituiti con doppi, e asserzioni su ci\`o che la schermata
        mostra e su ci\`o che parte al tocco, mai sul layout.
\end{itemize}

\section{Controlli statici}

\texttt{tsc} in modalit\`a strict, con \texttt{noUnusedLocals} e
\texttt{noUnusedParameters} attivi: il codice morto non passa. ESLint aggiunge
le regole di Expo (fra cui \texttt{react-hooks}, che ha fatto emergere due
effetti che scrivevano stato durante il render) e le quattro regole di
confine fra i livelli descritte nel capitolo sull'architettura.

\paragraph{Un difetto trovato dal linter appena installato.}
\texttt{expo/no-dynamic-env-var} ha segnalato che \texttt{env.ts} leggeva
\texttt{process.env[nome]} con un nome calcolato. La trasformazione Babel di
Expo sostituisce \texttt{process.env.EXPO\_PUBLIC\_X} \emph{testualmente} in
fase di build: con un accesso dinamico non c'\`e pi\`u nessun oggetto da
interrogare a runtime, quindi in una build installata quel ramo restituiva
sempre \texttt{undefined}. Funzionava tutto lo stesso solo perch\'e i valori
sono anche in \texttt{app.json}, ma le variabili registrate su EAS erano peso
morto. Ora ogni variabile \`e scritta per esteso.

\chapter{Build e distribuzione}

\section{Aggiornamento dello schema}

\texttt{supabase/schema.sql} \`e idempotente e va rieseguito dopo un
aggiornamento del codice: la revisione corrente aggiunge la funzione
\texttt{sposta\_occorrenze} (migrazione 0003), senza la quale riprogrammare una
data risponde ``operazione non riuscita'' --- esattamente come accadeva al
riordino degli allegati prima che \texttt{riordina\_allegati} fosse in
posizione. \`E l'unico passo manuale richiesto lato server, ed \`e anche
l'unico modo in cui questo progetto riesce a rompersi fra un aggiornamento del
codice e il primo tocco: il client parla di una funzione che il database non
ha ancora.

\section{Fase 0 — sviluppo (Expo Go)}
\texttt{npm start} sul Mac, scansione QR con Expo Go (stessa rete Wi-Fi).

\textbf{Limite di Expo Go}: gli allegati non funzionano. Il client OAuth
nativo di Google esige il redirect \texttt{<applicationId>:/oauthredirect},
che Expo Go non pu\`o registrare (l\`i lo schema diventa \texttt{exp://\dots},
che Google rifiuta). Tutto il resto --- login, ripassi, occorrenze, ricerca,
sincronizzazione --- funziona. Per lavorare sugli allegati serve una
development build (\texttt{expo-dev-client} \`e gi\`a fra le dipendenze).
Il tunnel ngrok si era rivelato inaffidabile e il pacchetto
\texttt{@expo/ngrok} \`e stato rimosso dalle dipendenze: era anche l'unica
vulnerabilit\`a segnalata da \texttt{npm audit} per cui non esisteva una
patch a monte.

\section{Sicurezza delle dipendenze}
\texttt{npm audit} riporta zero vulnerabilit\`a. Le quattro segnalate in
origine (\texttt{brace-expansion}, \texttt{postcss}, \texttt{tar},
\texttt{uuid}) erano tutte transitive del solo toolchain di build — Metro,
Jest, prebuild — e quindi non presenti nel bundle installato sul telefono.
Sono fissate con la sezione \texttt{overrides} di \texttt{package.json}
anzich\'e con \texttt{npm audit fix -{}-force}: quel comando avrebbe
aggiornato \texttt{expo} alla SDK~57 lasciando per\`o indietro
\texttt{react-native} e \texttt{babel-preset-expo}, producendo un insieme di
versioni incoerente e non funzionante. Per \texttt{brace-expansion} non
esiste una versione corretta nelle linee 1.x e 2.x, per cui l'override \`e
necessariamente globale. Ogni override va rimosso quando Expo aggiorna la
dipendenza a monte.

\section{Il peso dell'APK: da 109,3 a 27,7~MB}
\label{sec:peso}

Misurato release contro release, a codice applicativo invariato (in locale con
\texttt{./gradlew assembleRelease}; l'APK della build EAS \texttt{preview} che
ne \`e seguita pesa lo stesso). Non \`e un'ottimizzazione cercata: \`e la
scoperta che tre impostazioni davano per scontato il contrario di quello che
facevano.

\begin{center}
\small
\begin{tabular}{@{}lrr@{}}
\toprule
 & \textbf{prima} & \textbf{dopo} \\
\midrule
APK & 109,3~MB & \textbf{27,7~MB} \\
\texttt{lib/} native & 81~MB su 4 ABI & 16~MB su 1 ABI \\
\texttt{classes*.dex} & \textasciitilde39~MB & 9,1~MB \\
\texttt{libbarhopper\_v3.so} & 19,2~MB & assente \\
\bottomrule
\end{tabular}
\end{center}

Tre cause, in ordine di sorpresa.

\paragraph{Il dev client entrava in produzione.} \texttt{expo-dev-client}
finiva in \emph{ogni} build, release compresa, e si portava dietro ML Kit
barcode scanning (\texttt{libbarhopper\_v3.so}, 19,2~MB sulle quattro ABI) per
lo scanner di QR del menu sviluppatore, che questa app non ha. L'esclusione \`e
condizionale e non permanente --- il dev client serve davvero a
\texttt{expo run:android} e al profilo \texttt{development}, toglierlo per
sempre sarebbe rinunciarci --- e riguarda tutta la catena
(\texttt{dev-client}, \texttt{dev-launcher}, \texttt{dev-menu},
\texttt{dev-menu-interface}): da SDK~54 l'autolinking collega anche i moduli
transitivi, e ML Kit arriva da \texttt{dev-launcher}, non dal pacchetto
diretto. La prima versione dell'esclusione toglieva il pacchetto e lasciava
dentro tutto il resto. L'unica fonte che l'autolinking legge \`e
\texttt{expo.autolinking.exclude} di \texttt{package.json} --- non l'app config
--- quindi quella sezione la scrive uno script (\texttt{eas-build-post-install}),
che gira dopo l'installazione e prima di prebuild e Gradle.

\paragraph{Minify e shrink non erano attivi.} \texttt{build.gradle} le legge da
due property che, senza nessuno che le imposti, valgono \texttt{false}. Le
imposta ora \texttt{expo-build-properties} e non \texttt{gradle.properties} a
mano, perch\'e \texttt{android/} lo rigenera \texttt{prebuild} a ogni build EAS
e una modifica scritta a mano l\`i dentro non sopravvive. Il dex passa da
\textasciitilde39~MB a 9,1.

\paragraph{Quattro architetture in un APK che ne usa una.} Non era una svista
di configurazione: le librerie arrivano come AAR gi\`a compilati per tutte, e
\texttt{reactNativeArchitectures} vale solo per i moduli compilati da sorgente.
Due delle quattro (\texttt{x86}, \texttt{x86\_64}, 45~MB in due) esistono solo
per gli emulatori. Le toglie uno \texttt{splits.abi} aggiunto da un config
plugin (\texttt{plugins/withAbiRelease.js}). \texttt{ndk.abiFilters} era la
strada provata per prima e non toglie niente --- non si applica alle
\texttt{.so} gi\`a compilate che arrivano dalle dipendenze, che qui sono tutte:
verificato ricostruendo e ripesando.

\texttt{splits} peraltro non produce un APK pi\`u piccolo: ne produce
\emph{uno per ABI}, ciascuno con la sola architettura che serve, e
\texttt{expo run:android} sceglie da s\'e quello giusto per il dispositivo
collegato. Chi riceve che cosa:

\begin{center}
\small
\begin{tabular}{@{}p{3.4cm}p{3.1cm}p{8.0cm}@{}}
\toprule
\textbf{Profilo} & \textbf{ABI} & \textbf{Perch\'e} \\
\midrule
EAS \texttt{preview} & solo \texttt{arm64-v8a} &
  \`E l'unico profilo che produce un APK da installare a mano passandolo da un
  link: serve un file unico, e ogni telefono a 64 bit usa quella. Un
  dispositivo solo a 32 bit non lo installa --- compromesso accettabile per una
  distribuzione interna, non lo sarebbe per lo store. \\
\addlinespace
EAS \texttt{production} & tutte &
  \`E un app bundle: \`e il Play Store a consegnare a ogni dispositivo la sola
  architettura che gli serve, quindi tenerle tutte non pesa sul download di
  nessuno ed \`e ci\`o che mantiene installabile anche un telefono a 32 bit. \\
\addlinespace
build locali & tutte, una per APK &
  \`E il caso che restava scoperto, ed \`e il motivo per cui l'app installata
  da \texttt{expo run} pesava quanto quattro architetture invece che una. \\
\bottomrule
\end{tabular}
\end{center}

\texttt{npm run android} passa da \texttt{scripts/nativoPerApp.js}, che
rigenera il progetto nativo quando l'\texttt{applicationId} di
\texttt{android/} non corrisponde a quello del branch su cui ci si trova:
\texttt{expo run:android} lancia \texttt{prebuild} solo se la cartella
\emph{non esiste}, quindi dopo un cambio di branch produrrebbe un APK con
l'identificativo dell'app precedente --- che il dispositivo tratterebbe come un
aggiornamento di quella, sostituendola invece di affiancarla.
\texttt{scripts/releaseLocale.js} fa invece in locale le tre cose che su EAS
fanno il profilo e l'hook, e ripristina \texttt{package.json} a fine corsa,
riuscita o no: l'esclusione dall'autolinking \`e un dettaglio di build, e se
restasse sul disco finirebbe in un commit. Comandi, tabella e procedura per
rimisurare sono in \texttt{BUILD.md}.

\section{Fase 1 — app installate (vedi \texttt{BUILD.md})}
\begin{itemize}[nosep]
  \item \textbf{Android}: \texttt{eas build -p android --profile preview}
        $\rightarrow$ APK dal cloud Expo (gratis, \textasciitilde15
        build/mese), installazione diretta. \textbf{Fatto}: l'APK \`e stato
        costruito e installato, e pesa 27,7~MB (cap.~\ref{sec:peso}).
        Variabili d'ambiente registrate su EAS (il \texttt{.env} locale non
        viene caricato); il profilo disattiva l'upload automatico delle
        sourcemap a Sentry (\texttt{SENTRY\_DISABLE\_AUTO\_UPLOAD}), che senza
        un token d'organizzazione farebbe fallire la build.
  \item \textbf{iPad}: \texttt{npx expo prebuild -p ios} +
        \texttt{npx expo run:ios --device} con Apple ID gratuito (Personal
        Team). Limite: firma valida 7 giorni, rinnovo = rilanciare il
        comando con iPad collegato. Richiede Xcode e CocoaPods sul Mac.
  \item \textbf{Mac M2}: stessa app come ``My Mac (Designed for iPad)'' da
        Xcode. \`E il rischio 11.1 della specifica, in corso di validazione.
  \item Redirect OAuth per le build standalone: \texttt{ripassa://} va
        aggiunto agli URL consentiti in Supabase Auth.
\end{itemize}

\subsection*{Vincolo verificato: percorso senza spazi}
La build iOS nativa \`e stata provata end-to-end (Xcode 26.6). Fallisce se il
percorso del progetto contiene spazi — la cartella attuale
\texttt{Side hustle} rompe lo script Expo \texttt{get-app-config-ios.sh},
invocato senza virgolette (\texttt{bash: .../Documents/Side: No such file or
directory}). Spostando il progetto in un percorso senza spazi la stessa build
raggiunge \texttt{** BUILD SUCCEEDED **} e produce \texttt{Ripassa.app}
(binario universale arm64+x86\_64). \textbf{Il codice compila}: l'unico
ostacolo \`e lo spazio nel percorso. Android via cloud EAS non ne \`e
affetto. Dettagli e comandi in \texttt{BUILD.md}.

\chapter{Scostamenti dalla specifica iniziale}

\texttt{ripassa-app-spec.md} raccoglie i requisiti di luglio 2026 ed \`e
conservato come documento storico. Cinque punti sono stati superati
consapevolmente; elencarli \`e pi\`u onesto che riscrivere la specifica come
se ci avesse visto giusto dall'inizio.

\begin{center}
\small
\begin{tabular}{@{}p{3.5cm}p{5.2cm}p{5.9cm}@{}}
\toprule
\textbf{Specifica} & \textbf{Oggi} & \textbf{Perch\'e} \\
\midrule
\S2: nessuna dipendenza da Google Drive per lo storage &
  I binari degli allegati stanno sul Drive dell'utente &
  Il tetto di 1~GB di Supabase, gi\`a segnalato dalla specifica stessa come
  l'unico vincolo critico. Argomentazione completa nel cap.~\ref{sec:drive}. \\
\addlinespace
\S7.3: elimina dalla cache le righe con \texttt{cached\_at} fuori finestra &
  Elimina per appartenenza alla finestra &
  Il criterio originale cancellava e riscaricava file ancora validi solo
  perch\'e scaricati due giorni prima. Regressione coperta da test. \\
\addlinespace
\S4: tre livelli (Model, Controller, View) &
  Tre livelli pi\`u \texttt{config} trasversale &
  L'infrastruttura (client, SDK, storage cifrato) non appartiene al dominio e
  serve a tutti. Le dipendenze consentite sono ora verificate dal linter. \\
\addlinespace
\S8: Fase 0 su Expo Go come percorso completo &
  Gli allegati richiedono una build installata &
  Google esige un redirect OAuth derivato dall'\texttt{applicationId}, che
  Expo Go non pu\`o registrare. \\
\addlinespace
Nessuna notifica (fra le esclusioni di scopo) &
  Promemoria locali a ogni data programmata &
  L'esclusione riguardava le notifiche \emph{push} e resta valida: qui non c'\`e
  n\'e servizio push, n\'e token, n\'e costo, n\'e un server che sappia che cosa
  si ripassa. Un'app che ricorda le date senza ricordarle a nessuno era il
  difetto pi\`u evidente rimasto. Cap.~\ref{sec:notifiche}. \\
\bottomrule
\end{tabular}
\end{center}

Restano invece rispettati senza eccezioni: un solo account per tutti i
dispositivi, separazione MCV obbligatoria, budget, app nativa, nessuna
statistica, nessuna notifica \emph{push}.

\chapter{Metodo di lavoro e contributo personale}

Il progetto dichiara fin dalla specifica che il codice \`e generato da AI e
supervisionato, non scritto a mano. \`E una scelta di metodo, non una
scorciatoia, e come tutte le scelte di metodo va documentata --- altrimenti la
domanda ``e allora cosa hai fatto tu?'' resta senza risposta pur avendone una.

\section{Il ciclo effettivo}

\begin{enumerate}[nosep]
  \item \textbf{Specifica scritta a mano} prima di qualsiasi codice
        (\texttt{ripassa-app-spec.md}): obiettivo, decisioni vincolanti,
        modello dati, schermate, rischi da validare per primi. \`E il documento
        che rende generabile il resto.
  \item \textbf{Generazione} guidata, un'area funzionale alla volta.
  \item \textbf{Verifica}, con criteri di accettazione fissi: \texttt{npm run
        verifica} (lint, typecheck, test) deve passare; l'area toccata va
        provata sul dispositivo reale; le decisioni non ovvie devono risultare
        dal codice sotto forma di commento che spiega \emph{perch\'e}, non
        \emph{cosa}.
  \item \textbf{Correzione guidata} quando la verifica trova qualcosa --- che
        \`e il punto in cui si concentra il contributo personale, perch\'e
        quasi nessuno di questi difetti \`e visibile leggendo il codice.
\end{enumerate}

\section{Che cosa ha trovato la verifica}

Dieci esempi tracciabili nella storia dei commit. Sono la prova pi\`u concreta
di che cosa significhi ``supervisione'' qui.

\begin{center}
\scriptsize
\begin{tabular}{@{}p{1.5cm}p{5.4cm}p{7.7cm}@{}}
\toprule
\textbf{Commit} & \textbf{Problema} & \textbf{Come \`e emerso} \\
\midrule
\texttt{bf3a755} & La cache cancellava e riscaricava di continuo file ancora
  validi & Rilettura del criterio di eliminazione: la specifica stessa
  prescriveva la regola sbagliata (\texttt{cached\_at} invece di appartenenza) \\
\addlinespace
\texttt{255149e} & Sessione salvata in chiaro su disco; OAuth in implicit flow;
  RLS senza controllo sul genitore & Audit di sicurezza condotto rileggendo il
  codice con la domanda ``cosa vede chi ottiene il dispositivo?'' \\
\addlinespace
\texttt{2abbe0a} & I ripassi di oggi sparivano dalla sezione ``Da fare''
  appena passata l'ora & Uso reale dell'app: un difetto di prodotto, non di
  codice --- il confronto era per istante e doveva essere per giornata \\
\addlinespace
\texttt{5a8efbc} & Un \texttt{override} di sicurezza rompeva la build nativa
  Android & \texttt{expo run:android} fallito dopo \texttt{npm audit fix}:
  diagnosi del codegen di React Native e scelta di accettare la vulnerabilit\`a
  documentandola (\`e un DoS su pattern glob in build locale) \\
\addlinespace
\texttt{1581c99} & Google rifiutava lo schema \texttt{ripassa://} con
  \texttt{Error 400} & Prova sul dispositivo: i client nativi esigono
  \texttt{<applicationId>:/oauthredirect} \\
\addlinespace
\texttt{7b4cf7c} & Consenso concesso, ma l'app tornava al login senza errori &
  Riproduzione sul Samsung: Android uccide il processo mentre il browser \`e in
  primo piano, il redirect \emph{riavvia} l'app, e la promise che aspettava era
  morta col processo. Diagnosi impossibile da leggere nel codice \\
\addlinespace
\texttt{ef3c315} & Con token Drive revocato l'app diceva ``riavvia l'app'',
  l'unica azione inutile & Test dell'errore reale: \texttt{isAuthorized()}
  restava vero anche senza token utilizzabile \\
\addlinespace
\texttt{7c85747} & Correggere la prima data lasciava le altre quattro misurate
  dal giorno sbagliato & Uso reale: annotare oggi qualcosa studiato due giorni
  fa. Difetto di prodotto, invisibile nel codice --- ogni singola scrittura era
  corretta \\
\addlinespace
\texttt{68ec5f0} & L'APK pesava 109,3~MB: dev client e ML Kit in produzione,
  minify spento, quattro architetture su una usata & Pesatura dell'artefatto
  costruito, non lettura della configurazione. \texttt{ndk.abiFilters}, la
  strada che sembrava giusta, non toglieva niente: verificato ricostruendo e
  ripesando \\
\addlinespace
\texttt{f6b6496} & In modalit\`a aereo: rotella lunga e poi la schermata di
  login, l'unica che senza rete non serve a niente & Uso reale in assenza di
  rete: \texttt{getSession()} non \`e una lettura locale, e il suo
  \texttt{null} dice ``non lo so'' con le stesse parole di ``non c'\`e'' \\
\bottomrule
\end{tabular}
\end{center}

\section{Le quattro revisioni}

\begin{enumerate}[nosep]
  \item \textbf{MVP} (\texttt{14fbbb2}\dots\texttt{eb83eae}): funzionalit\`a,
        prima suite di test, audit di sicurezza, prima documentazione.
  \item \textbf{Affidabilit\`a e distribuzione}
        (\texttt{aebcea6}\dots\texttt{93787e9}): migrazione a Drive, errori
        traducibili, retry, connettivit\`a, crash reporting, e la lunga serie
        di correzioni sui due flussi OAuth.
  \item \textbf{Struttura e verifica} (\texttt{487b9ea} in poi):
        riorganizzazione per dominio dentro i livelli MCV a comportamento
        invariato, poi introduzione di linter, confini di architettura
        eseguibili, CI, test su repository e Controller, e riallineamento di
        questa documentazione al codice. In questa fase sono stati chiusi anche
        cinque difetti funzionali, fra cui il pi\`u grave dell'intero progetto:
        aprire la scheda di un ripasso prima che la lista fosse caricata e
        salvare cancellava titolo e note su tutti i dispositivi.
  \item \textbf{Uso quotidiano} (\texttt{1582832} in poi): la revisione delle
        cose che si scoprono solo usando l'app installata su un telefono vero.
        Promemoria locali, spostamento a cascata delle date, schermata di
        Profilo con Aiuto e segnalazione problemi, indicatore di ritento,
        apertura offline di sessione ed elenco, e il peso dell'APK sceso da
        109,3 a 27,7~MB. Sul fronte della verifica: le quattro schermate a
        copertura zero, i test della tabella delle regole d'errore e
        dell'impianto delle notifiche, con le soglie della \texttt{view}
        alzate da 29\,\% a 79\,\%.
\end{enumerate}

\section{Limite dichiarato del metodo}

Il codice generato \`e uniformemente buono in ci\`o che \`e locale
(nomi, struttura di una funzione, gestione di un caso) e sistematicamente
debole in ci\`o che dipende dal comportamento reale del sistema operativo,
dalla vita del processo e dall'uso quotidiano. Dieci dei dieci difetti
elencati sopra appartengono alla seconda categoria, e i tre aggiunti
nell'ultima revisione la allargano in una direzione che non avevo previsto: non
solo il sistema operativo e il ciclo di vita del processo, ma anche la
\emph{catena di build} --- un APK tre volte pi\`u pesante del necessario per
impostazioni che nessuno aveva scritto, e che quindi non si potevano leggere da
nessuna parte. \`E il motivo per cui la
verifica sul dispositivo non \`e sostituibile con i test, e per cui i test
servono comunque: intercettano la terza categoria, quella delle regressioni
introdotte correggendo le prime due.

\chapter{Rischi aperti}

\begin{enumerate}[nosep]
  \item Comportamento dei moduli nativi (fotocamera, picker) su macOS in
        modalit\`a ``Designed for iPad'' — da testare sul dispositivo reale.
  \item Rinnovo settimanale della firma iPad con Apple ID gratuito —
        fastidio accettabile o no? Decisione dopo l'uso.
  \item \textbf{Il viaggio attraverso il browser nei due flussi OAuth}:
        verificato a mano, non da test automatici (le ragioni sono nel
        cap.~\ref{sec:test}). Resta il rischio residuo pi\`u alto, perch\'e
        \`e anche l'area che storicamente ha richiesto pi\`u correzioni ---
        ridotto, non chiuso, dalle suite che ora coprono la parte
        deterministica di \texttt{driveAuth} e di \texttt{useAuth}.
  \item \textbf{La consegna effettiva dei promemoria}: i test verificano che
        l'app chieda il permesso, crei il canale e metta in coda la notifica
        giusta per l'istante giusto. Che il sistema operativo la mostri
        davvero, mesi dopo, con l'app mai aperta nel frattempo, non \`e una
        propriet\`a del codice: dipende da Doze, dalle ottimizzazioni della
        batteria di Samsung e da quanto l'utente sia stato invitato a
        ``ottimizzare'' l'app. \`E il tipo di guasto che si scopre non
        succedendo, cio\`e nel modo peggiore. Da verificare sull'uso reale, a
        distanza.
  \item \textbf{L'APK di \texttt{preview} \`e a sola \texttt{arm64-v8a}}
        (cap.~\ref{sec:peso}): su un dispositivo solo a 32 bit non si installa.
        \`E un compromesso accettato per la distribuzione a mano e non lo
        sarebbe per lo store --- il profilo \texttt{production} infatti tiene
        tutte le architetture. Il rischio \`e che la distinzione si perda:
        distribuire per errore l'APK di \texttt{preview} come se fosse la
        release \`e un guasto che si manifesta solo sul dispositivo di
        qualcun altro.
  \item \textbf{SHA-1 di release}: il client OAuth Android \`e registrato con
        l'impronta del keystore di debug. La prima build EAS di produzione
        user\`a un keystore diverso, e il login Drive fallirebbe solo su quella
        build. Procedura in \texttt{BUILD.md}, sottosezione
        \emph{«SHA-1: debug non \`e release»} (dentro la sezione~0,
        Prerequisiti, punto sul client OAuth Android). Allo stesso nodo
        appartiene il fatto che oggi esista un solo client OAuth, di tipo
        generico: \texttt{app.json} ne dichiarava tre, ma con lo stesso valore
        copiato tre volte, il che su Google Cloud non pu\`o essere corretto
        (il tipo di un client \`e fissato alla creazione). Le due chiavi
        ridondanti sono state tolte; \texttt{driveConfig.ts} le legge non
        appena i client per piattaforma esistono.
  \item Consumo dello spazio sul Drive personale con l'uso reale — non \`e
        pi\`u un tetto condiviso da 1~GB, ma resta da misurare.
\end{enumerate}

\paragraph{Chiusi nelle revisioni precedenti.} Il tetto di 1~GB dello storage
Supabase (non pi\`u applicabile: cap.~\ref{sec:drive}) e la compilazione iOS,
verificata end-to-end fino a \texttt{BUILD SUCCEEDED}.

\paragraph{Chiusi in questa revisione.} La build Android della Fase~1: l'APK
\`e costruito su EAS, installato e in uso, e pesa 27,7~MB
(cap.~\ref{sec:peso}) --- il che chiude anche il rischio, mai messo per
iscritto perch\'e nessuno lo aveva ancora misurato, di distribuire a mano
un'app da un centinaio di megabyte. Chiusa anche la copertura zero delle
quattro schermate, che rendeva ogni modifica alla \texttt{view} una modifica al
buio.

\end{document}
